\chapter{分子运动的统计描述}

\section{从状态概率到速度分布}
\label{th:mb}
宏观气体处于平衡时，单个分子的速度仍不断变化。统计描述用概率表示相同宏观条件下各微观状态的占据情况。若有 $N_i$ 个粒子占据第 $i$ 个单粒子状态，则 $p_i=N_i/N$，且 $\sum_i p_i=1$。在经典稀薄气体近似下，温度为 $T$ 的热平衡分布为
\begin{equation}
 p_i=A\ee^{-E_i/(\kb T)},\qquad A^{-1}=\sum_i\ee^{-E_i/(\kb T)}.
\end{equation}
若能量 $E_i$ 对应 $r_i$ 个不同状态，整个能级的概率应为 $r_iA\ee^{-E_i/(\kb T)}$。简并度计数与指数因子承担不同作用：前者计算可占据状态的数目，后者比较每个状态的权重。

这里的单粒子状态必须区分清楚。能量相同并不意味着状态相同；速度大小相同而方向不同，就是不同状态。经典分布还要求量子简并效应可忽略，粒子间相互作用对能量和计数的影响足够小。平动服从本讲分布，并不意味着转动、振动等内部自由度在任意温度都充分激发。

\subsection{速度分量的归一化}
\label{th:gaussian}
没有宏观流动与外场方向偏置时，平动能
\[
 E_k=\frac m2(v_x^2+v_y^2+v_z^2)
\]
使三维概率密度可分离为三个相同的一维因子。记 $\alpha=m/(2\kb T)$，一维密度为 $f_x(v_x)=A\ee^{-\alpha v_x^2}$。概率是密度乘速度区间，$\dd P=f_x(v_x)\dd v_x$，故 $f_x$ 的量纲是速率的倒数，而不是无量纲概率。

归一化需要 $\int_{-\infty}^{\infty}f_x\dd v_x=1$。先令 $I=\int_{-\infty}^{\infty}\ee^{-\alpha x^2}\dd x$，利用两个独立积分相乘，有
\begin{align}
 I^2&=\iint_{\mathbb R^2}\ee^{-\alpha(x^2+y^2)}\dd x\dd y\\
 &=\int_0^{2\pi}\dd\theta\int_0^\infty \ee^{-\alpha r^2}r\dd r\\
 &=2\pi\lim_{L\to\infty}
 \left[-\frac{\ee^{-\alpha r^2}}{2\alpha}\right]_0^L
 =\frac\pi\alpha.
\end{align}
因 $I>0$，得到 $I=\sqrt{\pi/\alpha}$，从而
\begin{equation}\label{eq:maxwellvector}
 f_x(v_x)=\sqrt{\frac{m}{2\pi\kb T}}\ee^{-mv_x^2/(2\kb T)},\quad
 f(\vct v)=\left(\frac{m}{2\pi\kb T}\right)^{3/2}\ee^{-mv^2/(2\kb T)}.
\end{equation}
三维归一化测度为 $\dd^3v=\dd v_x\dd v_y\dd v_z$。速度分量的独立性来自这里的平衡分布。非平衡气流若存在方向关联，联合密度就需要另行确定。

\Diagram{fig:gaussianarea}{高斯积分中直角坐标面积与极坐标面积的对应}{
 \draw[axis](-.4,0)--(4.3,0)node[right]{$x$};\draw[axis](0,-.4)--(0,3.3)node[above]{$y$};
 \draw[curve](0,0)circle(2.4);\draw[guide](0,0)circle(2.7);
 \fill[gray!20](0,0)--(35:2.7)arc(35:47:2.7)--(47:2.4)arc(47:35:2.4)--cycle;
 \draw(0,0)--(41:2.7);\node at(1.2,1.35){$r$};\node[anchor=west]at(2.9,1.9){$r\dd r\dd\theta$};
 \draw[->](.8,0)arc(0:41:.8);\node at(1.1,.35){$\theta$};
}
图\ref{fig:gaussianarea}中小扇环的径向宽度为 $\dd r$，切向长度为 $r\dd\theta$，两者相乘得到面积元。因子 $r$ 就来自切向长度随半径的增长。

\subsection{计算矩所需的积分}
\label{th:integrals}
定义半轴高斯矩
\[
 I_s(\alpha)=\int_0^\infty x^s\ee^{-\alpha x^2}\dd x,\qquad s>-1.
\]
偶函数的半轴积分为整轴的一半，故$I_0=\sqrt\pi/(2\sqrt\alpha)$。
对$I_1$取$u=\alpha x^2$，则$\dd u=2\alpha x\dd x$，端点仍为$0,\infty$，所以
\[
 I_1=\frac1{2\alpha}\int_0^\infty\ee^{-u}\dd u
 =\frac1{2\alpha}\lim_{U\to\infty}[-\ee^{-u}]_0^U
 =\frac1{2\alpha}.
\]
在积分号内对参数求导，$\partial_\alpha\ee^{-\alpha x^2}=-x^2\ee^{-\alpha x^2}$，因此
\[
 \frac{\dd I_s}{\dd\alpha}
 =-\int_0^\infty x^{s+2}\ee^{-\alpha x^2}\dd x=-I_{s+2}.
\]
递推式用分部积分建立。将$I_{s+2}$中的$x\ee^{-\alpha x^2}\dd x$
写成$-\dd(\ee^{-\alpha x^2})/(2\alpha)$，便有
\[
 \begin{aligned}
 I_{s+2}&=-\frac1{2\alpha}
 [x^{s+1}\ee^{-\alpha x^2}]_0^\infty
 +\frac{s+1}{2\alpha}\int_0^\infty x^s\ee^{-\alpha x^2}\dd x\\
 &=\frac{s+1}{2\alpha}I_s.
 \end{aligned}
\]
$s>-1$保证零端点的$x^{s+1}$趋零，无穷远端则由指数衰减消去边界项。
例如$I_4=3I_2/(2\alpha)=3I_0/(4\alpha^2)$、
$I_5=2I_3/\alpha=2I_1/\alpha^2$。需要频繁使用的结果为
\[
 \begin{array}{c|cccccc}
 s&0&1&2&3&4&5\\ \hline
 I_s&\dfrac{\sqrt\pi}{2\sqrt\alpha}&\dfrac1{2\alpha}&
 \dfrac{\sqrt\pi}{4\alpha^{3/2}}&\dfrac1{2\alpha^2}&
 \dfrac{3\sqrt\pi}{8\alpha^{5/2}}&\dfrac1{\alpha^3}
 \end{array}
\]
积分区间是半轴还是整轴，要在代入前确定。偶函数的整轴积分是半轴的两倍；奇函数的整轴积分为零，不能用半轴积分代替。

\subsection{平衡权重的来源与能量零点}
\label{th:reservoirweight}
Boltzmann 因子可由粒子与热库的能量交换推导。把两者组成孤立整体，忽略耦合所占的能量，总能量记为 $E_{\rm tot}$。粒子处于能量为 $E_i$ 的一个微观状态时，热库能量为 $E_{\rm tot}-E_i$；该状态的概率与热库在此能量下的可占据状态数 $\Omega_b(E_{\rm tot}-E_i)$ 成正比。用 $S_b=\kb\ln\Omega_b$ 表示热库熵，在参考能量附近展开，
\[
 S_b(E_b-E_i)=S_b(E_b)-\frac{E_i}{T}
 -\frac{E_i^2}{2C_bT^2}+\cdots.
\]
这里用到$\partial S_b/\partial E_b=1/T$。对它再用链式法则，并用$C_b=\partial E_b/\partial T$，
\[\frac{\partial^2S_b}{\partial E_b^2}
=\frac{\dd(T^{-1})}{\dd T}\frac{\partial T}{\partial E_b}
=-\frac1{T^2}\frac1{C_b}.\]
热库的热容$C_b$足够大时，粒子交换的有限能量对热库温度的改变可忽略，指数中保留一次项便得到$\ee^{-E_i/(\kb T)}$。这里的温度是共同平衡温度。若热库很小，二次项不可忽略，单一的正则指数权重便只是近似。

\Diagram{fig:reservoirweight}{粒子的状态权重由热库可占据状态数决定}{
 \node[draw,minimum width=2cm,minimum height=1.2cm](a)at(0,0){粒子 $E_i$};
 \node[draw,minimum width=3.3cm,minimum height=1.5cm](b)at(5,0){热库 $E_b-E_i$};
 \draw[<->](a)--(b)node[midway,above]{能量交换};
 \node at(2.5,-1.35){整体能量固定；热库温度近似不变};
}
图\ref{fig:reservoirweight}比较的是单个微观状态的权重。同一能级有若干状态时，将各状态概率相加便出现简并度因子。连续能量情形相应地乘以态密度，因此能量分布的峰值由态密度与指数权重共同决定。

设单粒子配分和为 $Z_1=\sum_i r_i\ee^{-\beta E_i}$，其中 $\beta=1/(\kb T)$。若所有能级同时平移常数 $E_\ast$，则
\[
 Z_1'= \ee^{-\beta E_\ast}Z_1,\qquad
 \frac{r_i\ee^{-\beta(E_i+E_\ast)}}{Z_1'}
 =\frac{r_i\ee^{-\beta E_i}}{Z_1}.
\]
占有概率不变，平均能量只增加 $E_\ast$。粒子数固定时，能量差、做功和热容都不受这一平移影响；粒子数变化的开口系统则须在储能和粒子携带能量两处采用同一零点。

\subsection{由各向同性与分量独立性确定高斯形状}
各向同性规定分布只能依赖 $v_x^2+v_y^2+v_z^2$；分量独立性又规定联合密度可写成三个一维因子的乘积。这两个条件共同限制了函数形状。写成
\[
 f(\vct v)=\phi(v_x)\phi(v_y)\phi(v_z)
 =\Psi(v_x^2+v_y^2+v_z^2).
\]
在密度为正、二次连续可微的区域取对数，记 $L(s)=\ln\Psi(s)$。对 $v_x$ 求导，再对 $v_y$ 求导，
\[
 \frac{\partial^2\ln f}{\partial v_y\partial v_x}
 =4v_xv_yL''(v^2)=0.
\]
左端为零是因为乘积的对数中没有同时含$v_x,v_y$的项。任取非零的$v_x,v_y$可知$L''(s)=0$。两次积分分别给$L'(s)=-\alpha$、$L(s)=a-\alpha s$，其中$a,\alpha$为常数。归一化要求$\alpha>0$，于是得到高斯形状。由其二阶矩可得$m\langle v_x^2\rangle=m/(2\alpha)$；与平衡态的热能尺度$\kb T$对应，便有$\alpha=m/(2\kb T)$。

各向同性单独并不保证高斯分布。例如，速率固定而方向均匀的粒子群也是各向同性的，但其三个速度分量受 $v_x^2+v_y^2+v_z^2=v_0^2$ 约束，彼此不独立。后文固定速率背景的碰撞题正属于这种分布，不能把它的某一分量再当成 Maxwell 高斯分量。

\subsection{高斯积分的缩放、递推与边界}
已有的 $I_s$ 表可以由一个无量纲积分统一使用。令 $y=\sqrt\alpha\,x$，则
\[
 I_s(\alpha)=\alpha^{-(s+1)/2}
 \int_0^\infty y^s\ee^{-y^2}\dd y.
\]
幂次 $\alpha^{-(s+1)/2}$ 同时包含变量幂的缩放和微分 $\dd x$ 的缩放。它可先用来判断答案的温度、质量依赖，再计算常数系数。对连续参数求导时，积分区间须与该参数无关，并要求导数的积分收敛；本讲的 $\alpha>0$ 高斯积分满足这些条件。

若下限为非零阈值 $a$，则定义
$J_s(a)=\int_a^\infty x^s\ee^{-\alpha x^2}\dd x$。分部积分保留下限项，得到
\[
\begin{aligned}
 J_{s+2}(a)&=-\frac1{2\alpha}[x^{s+1}e^{-\alpha x^2}]_a^\infty
+\frac{s+1}{2\alpha}\int_a^\infty x^se^{-\alpha x^2}\dd x\\
&=\frac{a^{s+1}\ee^{-\alpha a^2}}{2\alpha}
 +\frac{s+1}{2\alpha}J_s(a).
\end{aligned}
\]
取 $a=0$ 才回到 $I_s$ 的递推。阈值问题出现的“指数乘多项式”，正是反复保留下限项的结果；半轴归一化积分中的边界项为零，并不意味着截断积分的边界项也为零。


\begin{example}[两个能级的占有与热容]\label{ex:twolevels}
独立粒子只有能量 $0$ 和 $\varepsilon$ 两个能级，简并度分别为 $1$、$r$。粒子数 $N$ 固定，求激发态占比、内能和定容热容，并讨论高、低温极限。
\Diagram{fig:twolevelq}{两个能级及其简并度}{
 \draw[axis](0,0)--(0,2.4)node[above]{$E$};\draw[curve](.5,.3)--(3.5,.3)node[right]{$0\;(1)$};
 \foreach \x in{.5,1.55,2.6}{\draw[curve](\x,1.65)--({\x+.8},1.65);}\node[right]at(3.5,1.65){$\varepsilon\;(r)$};
 \draw[->](2,.5)--(2,1.35)node[midway,right]{$T$};
}
\end{example}
\begin{solution}
图\ref{fig:twolevelq}中的上能级包含 $r$ 个状态。把各状态权重相加，单粒子归一化因子为 $Z=1+r\ee^{-x}$，其中 $x=\varepsilon/(\kb T)$。于是
\[
 P_1=\frac{r\ee^{-x}}{1+r\ee^{-x}},\qquad U=N\varepsilon P_1.
\]
由于$x=\varepsilon T^{-1}/\kb$，有$\dd x/\dd T=-\varepsilon/(\kb T^2)=-x/T$。
占有率的商式求导逐项给出
\[
 \frac{\dd P_1}{\dd x}
 =\frac{(-r\ee^{-x})(1+r\ee^{-x})-(r\ee^{-x})(-r\ee^{-x})}
 {(1+r\ee^{-x})^2}
 =-\frac{r\ee^{-x}}{(1+r\ee^{-x})^2}.
\]
再对$U=N\varepsilon P_1(x(T))$使用链式法则，
\[
 C_V=N\varepsilon\frac{\dd P_1}{\dd x}\frac{\dd x}{\dd T}
 =N\kb x^2\frac{r\ee^{-x}}{(1+r\ee^{-x})^2}.
\]
热容的温度导数是在粒子数不变的条件下取的。还可把结果中的分式写成 $P_1(1-P_1)$：占有率接近0或1时，粒子很少改变能级，内能随温度的变化便很小。对单粒子有 $\langle E^2\rangle=\varepsilon^2P_1$，所以
\[
 C_V=\frac{N}{\kb T^2}
 \bigl(\langle E^2\rangle-\langle E\rangle^2\bigr).
\]
这里的方差指一个粒子的能量涨落。$N$ 个独立粒子的总方差为它的 $N$ 倍，与上面直接求导相符。高温极限中各微观状态等概率，上能级占比仍由其简并度决定。

\Diagram{fig:twolevela}{能量与温度的两个极限}{
 \draw[axis](0,0)--(5,0)node[right]{$T$};\draw[axis](0,0)--(0,2.3)node[above]{$U/N$};
 \draw[curve]plot[domain=.12:4.5,samples=60](\x,{1.8*exp(-1/\x)/(1+exp(-1/\x))});
 \draw[guide](0,.9)--(4.6,.9)node[right]{$r\varepsilon/(1+r)$};\node at(3.5,1.7){示意取 $r=1$};
}
图\ref{fig:twolevela}显示有限能级系统的内能有上界。低温 $x\gg1$ 时激发被冻结，$C_V\sim N\kb r x^2\ee^{-x}$；高温 $x\ll1$ 时各微观状态近乎等概率，$P_1\to r/(1+r)$。此时内能接近饱和值，继续升温引起的内能增量趋于零，热容也趋于零。这与具有无界二次型能量的经典自由度不同。
\end{solution}

\section{速率、方向与抽样方式}
\label{th:speed}
速度矢量落入速度空间体积元的概率为 $f(\vct v)\dd^3v$。若只问速率在 $v$ 与 $v+\dd v$ 之间，就要把所有方向相加。各向同性时，球壳体积为 $4\pi v^2\dd v$，故
\begin{equation}\label{eq:speeddensity}
 g(v)=4\pi\left(\frac{m}{2\pi\kb T}\right)^{3/2}v^2\ee^{-mv^2/(2\kb T)},\qquad v\ge0.
\end{equation}
\Diagram{fig:velocityshell}{速度空间中厚度为\(\dd v\)的球壳}{
 \draw[axis](-2.1,0)--(2.4,0)node[right]{$v_x$};\draw[axis](0,-1.8)--(0,2.1)node[above]{$v_z$};
 \draw[axis](-1.7,-1.2)--(1.7,1.2)node[right]{$v_y$};
 \draw[curve](0,0)circle(1.45);\draw[guide](0,0)circle(1.62);
 \draw[curve](0,0)ellipse(1.45 and .4);\draw[guide](0,0)ellipse(1.62 and .45);
 \draw[->](0,0)--(35:1.45)node[midway,above]{$v$};
 \draw[<->](315:1.45)--(315:1.62);\node[right]at(1.5,-1.2){$\dd v$};
}
$f(\vct v)$ 在原点最大，而 $g(v)$ 在 $v=0$ 为零。这两个结论并不矛盾：球壳因子 $v^2$ 把不同方向的状态数计入速率概率。图\ref{fig:velocityshell}中的球壳是概率求和的几何区域，并非分子在真实空间的轨道。

\subsection{三个特征速率与一般矩}
\label{th:moments}
令 $B=4\pi[m/(2\pi\kb T)]^{3/2}$，则
\[
 \langle v^s\rangle=B I_{s+2}(\alpha)
 =\frac2{\sqrt\pi}\Gamma\!\left(\frac{s+3}{2}\right)
 \left(\frac{2\kb T}{m}\right)^{s/2}.
\]
这里的一般矩由$u=\alpha v^2$得到：$\dd v=\dd u/(2\sqrt{\alpha u})$，因而
\[
 I_{s+2}=\frac1{2\alpha^{(s+3)/2}}
 \int_0^\infty u^{(s+1)/2}\ee^{-u}\dd u
 =\frac{\Gamma((s+3)/2)}{2\alpha^{(s+3)/2}},\qquad s>-3.
\]
最后的积分定义了$\Gamma$函数。整数矩也可直接使用前节递推。平均速率由$I_3$给出：
\[
 \bar v=B\frac1{2\alpha^2}=\sqrt{\frac{8\kb T}{\pi m}}.
\]
均方速率由 $I_4$ 给出，$\langle v^2\rangle=3\kb T/m$，故 $v_{\rm rms}=\sqrt{3\kb T/m}$。最概然速率需对密度 $g(v)$ 求极值：
\[
 g'(v)=B[2v\ee^{-\alpha v^2}+v^2(-2\alpha v)\ee^{-\alpha v^2}]
 =2Bv\ee^{-\alpha v^2}(1-\alpha v^2)=0,
 \qquad v_p=\sqrt{\frac{2\kb T}{m}}.
\]
除 $v=0$ 端点外，只有一个峰值。因此 $v_p<\bar v<v_{\rm rms}$。此外
\[
 \langle v^3\rangle=\sqrt{\frac{128}\pi}\left(\frac{\kb T}{m}\right)^{3/2},\quad
 \langle v^4\rangle=15\left(\frac{\kb T}{m}\right)^2.
\]
等分定理在这里可直接验证：$\langle v_x^2\rangle=\kb T/m$，每个平动二次项的平均能量为 $\kb T/2$。能量由二阶矩决定，碰撞通量由一阶矩决定，不能用同一个“平均速度”处理所有问题。

\begin{example}[能量密度的最概然值]
求平动能 $E=mv^2/2$ 的概率密度 $h(E)$、平均值与最概然值。为什么把 $v_p$ 代入 $mv^2/2$，不等于能量密度的峰值位置？
\Diagram{fig:energydensityq}{速率区间映射为不等宽的能量区间}{
 \draw[axis](0,1.7)--(5,1.7)node[right]{$v$};\draw[axis](0,0)--(5,0)node[right]{$E$};
 \draw[line width=2pt](1,1.7)--(1.5,1.7);\draw[line width=2pt](1,0)--(1.8,0);
 \draw[->](1.25,1.4)--(1.4,.3)node[midway,right]{$\dd E=mv\dd v$};
}
\end{example}
\begin{solution}
同一批粒子的概率不因变量替换而变，图\ref{fig:energydensityq}中两段对应 $h(E)\dd E=g(v)\dd v$。由 $v=\sqrt{2E/m}$、$\dd v/\dd E=1/\sqrt{2mE}$，得到
\[
 h(E)=\frac2{\sqrt\pi}\frac{\sqrt E}{(\kb T)^{3/2}}\ee^{-E/(\kb T)},\quad E\ge0.
\]
平均值可直接由二阶速率矩得到 $\langle E\rangle=3\kb T/2$。对密度取对数后求导，
\[
 \frac{\dd\ln h}{\dd E}=\frac1{2E}-\frac1{\kb T}=0
 \quad\Rightarrow\quad E_p=\frac12\kb T.
\]
变量替换中的幂次可以逐项追踪：原密度里的 $v^2$ 变为与 $E$ 成正比，雅可比里的 $1/v$ 又给出 $E^{-1/2}$，相乘才得到 $E^{1/2}$。密度的量纲也从 $1/v$ 变为 $1/E$。例如一小段能量区间的概率近似为 $h(E)\Delta E$；比较两个区间的粒子数时，既要比较密度，又要比较区间宽度。以速率等宽和以能量等宽划分样本，所比较的本来就是不同的区间。

\Diagram{fig:energydensitya}{能量概率密度的峰值与平均值}{
 \begin{scope}[yscale=3]
 \draw[axis](0,0)--(5,0)node[right]{$E/(\kb T)$};\draw[axis](0,0)--(0,.85)node[above]{$\kb T h$};
 \draw[curve]plot[domain=0:4.5,samples=90](\x,{2/sqrt(pi)*sqrt(\x)*exp(-\x)});
 \draw[guide](.5,0)--(.5,.484)node[above left]{$E_p$};\draw[guide](1.5,0)--(1.5,.308)node[above right]{$\langle E\rangle$};
 \end{scope}
}
图\ref{fig:energydensitya}的横轴为能量。用 $mv_p^2/2$ 得到 $\kb T$，它只表示速率密度峰值所对应的能量，不是 $h(E)$ 的峰值。概率密度的极大位置会受变量变换中的雅可比影响；具有直接不变意义的是某一区间的积分概率。
\end{solution}

\begin{example}[速度分量的阈值与尾概率]
平衡气体温度为 $T$。求 $v_x>u$ 的粒子比例及其中的条件平均 $\langle v_x\mid v_x>u\rangle$，其中 $u\ge0$。
\Diagram{fig:tailq}{一维速度分布中的阈值区域}{
 \draw[axis](-3,0)--(3,0)node[right]{$v_x$};\draw[axis](0,0)--(0,2.3)node[above]{$f_x$};
 \fill[gray!20](1,0)--plot[domain=1:2.7](\x,{1.8*exp(-\x*\x)})--(2.7,0)--cycle;
 \draw[curve]plot[domain=-2.7:2.7,samples=70](\x,{1.8*exp(-\x*\x)});\node[below]at(1,0){$u$};
}
\end{example}
\begin{solution}
按图\ref{fig:tailq}只积分右侧阴影，不再乘二。设 $z=\sqrt\alpha u$，定义
\[
 \erfc z=(2/\sqrt\pi)\int_z^\infty\ee^{-t^2}\dd t,
\]
则
\[
 P(v_x>u)=\int_u^\infty\sqrt{\frac\alpha\pi}\ee^{-\alpha v_x^2}\dd v_x
=\frac1{\sqrt\pi}\int_z^\infty\ee^{-t^2}\dd t
=\frac12\erfc z.
\]
其中$t=\sqrt\alpha v_x$，$\dd v_x=\dd t/\sqrt\alpha$，下限$u$变为$z$，上限仍为无穷大。
条件平均的分子为 $\int_u^\infty v_x\sqrt{\alpha/\pi}\ee^{-\alpha v_x^2}\dd v_x$。利用 $\dd(\ee^{-\alpha v_x^2})=-2\alpha v_x\ee^{-\alpha v_x^2}\dd v_x$，有
\[
 \int_u^\infty v_x f_x\dd v_x
 =\sqrt{\frac\alpha\pi}\lim_{L\to\infty}
 \left[-\frac{\ee^{-\alpha v_x^2}}{2\alpha}\right]_u^L
 =\frac{\ee^{-z^2}}{2\sqrt{\pi\alpha}},\qquad
 \langle v_x\mid v_x>u\rangle=
 \frac{\ee^{-z^2}}{\sqrt{\pi\alpha}\,\erfc z}.
\]
条件抽样可直接写成归一化密度
\[
 q(v_x)=\frac{f_x(v_x)}{P(v_x>u)},\qquad v_x>u,
\]
阈值之外密度为零。分子积分只给出原总体中这部分粒子贡献的速度一阶矩，除以尾概率才是子样本的平均值。对 $u=0$，原总体一半粒子被选中；对 $u>0$，所选比例进一步减小。由于保留的每个速度都大于 $u$，所得条件平均也应大于 $u$，这与把一个未经归一化的一阶矩误当平均值的结果有明显区别。

\Diagram{fig:taila}{阈值右侧重新归一化后的条件密度}{
 \draw[axis](-.3,0)--(4.5,0)node[right]{$v_x$};\draw[axis](0,0)--(0,3.3)node[above]{$q$};
 \fill[gray!15](1,0)--plot[domain=1:3.5,samples=70](\x,{7.1739*exp(-\x*\x)})--(3.5,0)--cycle;
 \draw[curve](1,0)--(1,{7.1739*exp(-1)});\draw[curve]plot[domain=1:3.5,samples=70](\x,{7.1739*exp(-\x*\x)});
 \node[below]at(1,0){$u$};\node at(2.8,1.8){面积为1};\node at(2.3,-.7){$q=f_x/P(v_x>u)$};
}
$u=0$ 时结果为 $\sqrt{2\kb T/(\pi m)}$。图\ref{fig:taila}强调，这是从某一时刻的气体中筛选速度分量；若按单位时间撞壁的粒子抽样，还要额外乘碰撞速率，不能沿用这个条件密度。
\end{solution}

\subsection{同一粒子群的不同概率测度}
\label{th:samplingmeasure}
概率密度总是相对于指定的积分变量定义。若一个单调变量变换为 $y=F(x)$，相同小区间中的概率应相等，
\[
 q(y)\dd y=p(x)\dd x,\qquad
 q(y)=p(x(y))\left|\frac{\dd x}{\dd y}\right|.
\]
密度的峰值会随变量变换而改变，区间概率和由同一事件定义的期望值则保持一致。例如直接计算 $\langle mv^2/2\rangle$，与先变成能量密度再计算 $\langle E\rangle$，必得相同结果。两种算法若不一致，应检查微分因子和积分范围。

\Diagram{fig:measuremap}{变量变换保持概率而改变区间宽度}{
 \draw[axis](0,1.4)--(6,1.4)node[right]{$x$};
 \draw[axis](0,0)--(6,0)node[right]{$y$};
 \draw[line width=2pt](1,1.4)--(1.8,1.4);
 \draw[line width=2pt](2.3,0)--(4.3,0);
 \draw[->](1.4,1.15)--(3.3,.25);
 \node[above]at(1.4,1.55){$\dd x$};
 \node[below]at(3.3,-.15){$\dd y$};
 \node[right]at(4,.8){$p(x)\dd x=q(y)\dd y$};
}
图\ref{fig:measuremap}中，两条粗线代表同一批事件在两个变量下占据的区间。速度矢量变为速率时，还须合并所有方向上的事件，球壳面积便来自这一步求和。这与单纯将一个一维变量改写成另一个变量有所不同。

另一种变化来自抽样方式。若一个粒子处于状态 $x$ 时，每单位时间被记录的机会正比于 $w(x)\ge0$，原有归一化密度为 $p(x)$，则记录事件的密度为
\[
 p_{\rm event}(x)=\frac{w(x)p(x)}{\langle w\rangle},
 \qquad
 \langle A\rangle_{\rm event}
 =\frac{\langle Aw\rangle}{\langle w\rangle}.
\]
分母是总记录率，用来把加权结果重新归一化。两次平均一般不同；差值可写成
\[
 \langle A\rangle_{\rm event}-\langle A\rangle
 =\frac{\langle Aw\rangle-\langle A\rangle\langle w\rangle}
 {\langle w\rangle}.
\]
因此事件样本中某个量偏大还是偏小，取决于它与记录权重的关联。快分子更常穿孔，所以穿孔样本的能量偏大；在一段束流中作瞬时空间计数，慢分子停留更久，权重又不同。这两种偏置都可由同一归一化式推导。

\subsection{条件筛选与通量加权的次序}
探测器若只接受某一区域 $\mathcal A$，先用指示函数 $\mathbf1_{\mathcal A}$ 表示筛选，则被记录事件的概率为
\[
 \dd P_{\rm det}=
 \frac{\mathbf1_{\mathcal A}(x)w(x)p(x)\dd x}
 {\int_{\mathcal A}w(x)p(x)\dd x}.
\]
分子中的筛选因子和权重可交换乘法次序，分母则始终取最终被接受事件的总记录率。体内满足阈值的粒子比例按粒子数计算，通量比例还包含各粒子的穿面机会，两者应分别积分。

各向同性小孔泄流中，速率因子 $vg(v)$ 与角度因子 $\cos\theta\,\dd\Omega$ 分离。单独限制接受圆锥不会改变速率分布；单独限制速率也不会改变归一化的角分布。这一结论以孔的传输几何与速率无关为条件。带电粒子经偏转场、长孔中的碰壁损失或随速率变化的探测效率，都可能使速率和方向在记录样本中重新关联。

\subsection{累积分布、典型宽度与尾部估计}
分量高斯密度的自然宽度是 $s=\sqrt{\kb T/m}$，其分布写成
\[
 f_x(v_x)=\frac1{\sqrt{2\pi}s}
 \exp\!\left(-\frac{v_x^2}{2s^2}\right).
\]
两个对称阈值之间的概率是
$P(|v_x|<a)=\erf(a/(\sqrt2s))$，阈值之外则为相应的 $\erfc$。所谓“宽度约为热速率”，只表示密度发生显著变化的尺度；精确的区间概率仍须积分。

对于远尾 $z\gg1$，令 $K(z)=\int_z^\infty\ee^{-t^2}\dd t$，用 $\ee^{-t^2}\dd t=-\dd(\ee^{-t^2})/(2t)$ 分部积分，
\[
 K(z)=\frac{\ee^{-z^2}}{2z}
 -\frac12\int_z^\infty\frac{\ee^{-t^2}}{t^2}\dd t.
\]
被减去的积分为正，因此首项给上界。再利用 $t^{-2}\le z^{-2}$，
\[
 \frac{\ee^{-z^2}}{2z+z^{-1}}
 \le K(z)\le\frac{\ee^{-z^2}}{2z}.
\]
两界之比趋于1，故
$\erfc z\sim\ee^{-z^2}/(\sqrt\pi z)$。指数给主要压低作用，前面的 $1/z$ 来自尾区有效宽度。对较高阈值估算泄流、反应或探测计数时，这一前因子也应保留。

\subsection{速率谱、能量谱与仪器记录的转换}
实际记录的横坐标可以是速度、能量或飞行时间。先把记录事件的密度写成 $p_v(v)$，再换横坐标，才能保持同一批样本。设测量路程固定为 $\ell$，到达时间 $t=\ell/v$，则
\[
 p_t(t)=p_v(\ell/t)\frac{\ell}{t^2},\qquad
 p_E(E)=p_v\!\left(\sqrt{\frac{2E}{m}}\right)
 \frac1{\sqrt{2mE}}.
\]
两个雅可比分别来自 $\left|\dd v/\dd t\right|$ 与
$\left|\dd v/\dd E\right|$。例如，等宽时间窗口并不对应等宽速度窗口：早到粒子的相邻时间窗口覆盖更大的速度区间。图上计数峰的位置因此随记录变量变化，测温时应拟合整个已归一化的记录模型，而不宜只把某个峰当作最概然速率。

若仪器记录效率为 $\eta(v)$，实际事件密度还需乘此因子，
$p_{\rm rec}(v)=\eta(v)p_v(v)/\langle\eta\rangle$。
反推出原谱时，可以在效率非零的区域除去该权重并重新归一化；效率为零的区间没有记录信息，单凭已测区间无法恢复其概率。有限飞行路径分散和有限时间分辨率则引入不同初始条件的叠加，属于在变量变换之外的另一层平均。


\section{相对运动、碰撞与自由程}
\label{th:relative}
取同温、互不相关的两个粒子，质量为 $m_1,m_2$，令 $M_t=m_1+m_2$、$\mu=m_1m_2/M_t$。相对速度与质心速度定义为
\[
 \vct u=\vct v_1-\vct v_2,\qquad
 \vct V_c=\frac{m_1\vct v_1+m_2\vct v_2}{M_t}.
\]
反解为 $\vct v_1=\vct V_c+(m_2/M_t)\vct u$、$\vct v_2=\vct V_c-(m_1/M_t)\vct u$。沿任一坐标方向的变换行列式为
\[
 \det\frac{\partial(v_1,v_2)}{\partial(u,V_c)}
 =\begin{vmatrix}m_2/M_t&1\\-m_1/M_t&1\end{vmatrix}=1.
\]
因而六维测度也满足 $\dd^3v_1\dd^3v_2=\dd^3u\dd^3V_c$。概率变量变换中使用的是行列式绝对值。一般非正交变换把小矩形变成平行四边形，面积不是随意把两个微分相乘所得的结果。

把反解代入动能，交叉项的系数为$2[m_1m_2/M_t-m_2m_1/M_t]=0$，相对速度平方的系数为$[m_1m_2^2+m_2m_1^2]/M_t^2=\mu$，所以
\[
 m_1v_1^2+m_2v_2^2=M_tV_c^2+\mu u^2.
\]
联合密度的指数因而分成$\exp[-M_tV_c^2/(2\kb T)]$与$\exp[-\mu u^2/(2\kb T)]$两因子。对质心速度的积分使用式\eqref{eq:maxwellvector}已证的三维高斯归一化，其值为$(2\pi\kb T/M_t)^{3/2}$。与原联合密度的系数相乘后，利用$m_1m_2/M_t=\mu$，得到
\begin{equation}\label{eq:relative}
 f_u(\vct u)=\left(\frac\mu{2\pi\kb T}\right)^{3/2}
 \ee^{-\mu u^2/(2\kb T)},\qquad
 \bar u=\sqrt{\frac{8\kb T}{\pi\mu}},\quad
 \langle u^2\rangle=\frac{3\kb T}\mu.
\end{equation}
这里 $\langle\vct u\rangle=\vct0$，而 $\bar u=\langle|\vct u|\rangle>0$。矢量平均与速率平均不能混写。

\begin{example}[不同分子之间的碰撞速率]\label{ex:relative}
温度 $T$ 下，气体组分1的质量为 $m_1$、数密度为 $n_1$，组分2为 $m_2,n_2$。硬球半径分别为 $r_1,r_2$。求每个组分1粒子与组分2粒子的平均碰撞频率，以及单位体积、单位时间的异种碰撞次数。
\Diagram{fig:relativeq}{实验室系的两速度与粒子2参照系中的相对运动}{
 \node at(1.1,1.7){实验室系};
 \node[particle,label=below:$1$](a)at(0,0){};\node[particle,label=below:$2$](b)at(2.5,0){};
 \draw[axis](a)--(1.1,.5)node[above]{$\vct v_1$};\draw[axis](b)--(3.3,.5)node[right]{$\vct v_2$};
 \begin{scope}[xshift=5.4cm]
 \node at(1.2,1.7){以粒子2为参照};\node[particle,label=below:$2$]at(2.2,0){};
 \draw[guide](2.2,0)circle(.7);\node[particle,label=below:$1$]at(0,.4){};
 \draw[axis](0,.4)--(1.25,.4)node[midway,above]{$\vct u$};
 \draw[<->](2.2,0)--(2.2,.7);\node[right]at(2.3,.8){$r_1+r_2$};
 \end{scope}
}
\end{example}
\begin{solution}
图\ref{fig:relativeq}中两球接触时中心距为 $r_1+r_2$。把粒子2的中心固定，另一中心以相对速度 $\vct u$ 运动，对应碰撞截面
\[
 \sigma_{12}=\pi(r_1+r_2)^2.
\]
在一小段时间 $\dd t$ 内，相对运动扫过体积为 $\sigma_{12}u\dd t$。先对速度平均，再乘靶粒子数密度，得到单个粒子的平均频率
\[
 z_{12}=n_2\sigma_{12}\bar u
 =n_2\pi(r_1+r_2)^2\sqrt{\frac{8\kb T}{\pi\mu}}.
\]
式中需要的是两粒子中心之间的相对扫掠长度。若直接把组分1的平均速率代入，会把组分2也在运动这一事实漏掉。对独立速度分量有
\[
 \langle u_x^2\rangle=\langle v_{1x}^2\rangle+
 \langle v_{2x}^2\rangle
 =\kb T\left(\frac1{m_1}+\frac1{m_2}\right).
\]
括号正好是 $1/\mu$，与前文相对速度分布的宽度一致。只有求出这一相对分布的一阶速率矩后，才能与截面、靶密度相乘。截面按中心接触条件定义，半径已经包含两球半径之和。

\Diagram{fig:relativea}{相对运动扫掠的碰撞圆柱}{
 \draw[curve](0,-.6)--(4,-.6)arc(-90:90:.2 and .6)--(0,.6);
 \draw[curve](0,0)ellipse(.2 and .6);\draw[axis](.2,0)--(3.8,0)node[midway,above]{$u\dd t$};
 \node at(2,-1.1){横截面 $\sigma_{12}$};
}
图\ref{fig:relativea}的每个圆柱以一个组分1粒子为起点，故体积碰撞率为 $R_{12}=n_1z_{12}$。每一对异种粒子按“1碰2”计数一次，无须再除以二。若统计同一种气体的所有碰撞，同一事件会从两个参与者各计一次，结果应除以二。
\end{solution}

\subsection{平均自由程与无碰撞概率}
\label{th:collisions}
同种硬球半径为 $r$ 时，$\sigma=4\pi r^2$，$\bar u=\sqrt2\bar v$。用平均率 $z=n\sigma\bar u$ 近似碰撞过程，若 $P(t)$ 表示从一次碰撞后持续到 $t$ 仍未再次碰撞的概率，则
\[
 P(t+\dd t)=P(t)(1-z\dd t),\qquad
 \dot P=-zP,\qquad P(0)=1.
\]
分离变量并用初始概率归一化，
\[\int_1^{P(t)}\frac{\dd p}{p}=-z\int_0^t\dd s,
\qquad[\ln p]_1^{P(t)}=-zt,\qquad P(t)=\ee^{-zt}.\]
等待时间密度为$w(t)=-\dot P=z\ee^{-zt}$。其归一化与平均时间分别为
\[\begin{aligned}
\int_0^\infty w(t)\dd t&=[-e^{-zt}]_0^\infty=1,\\
\tau=\int_0^\infty tw(t)\dd t
&=\left[-\left(t+\frac1z\right)e^{-zt}\right]_0^\infty=\frac1z.
\end{aligned}\]
第二个原函数可由$\int t e^{-zt}\dd t$分部积分得到。平均自由程采用分子运动论的平均定义
\begin{equation}
 \lambda=\bar v\tau=\frac{\bar v}{n\sigma\bar u}
=\frac{\bar v}{n\sigma\sqrt2\bar v}=\frac1{\sqrt2n\sigma}.
\end{equation}
严格说，给定速度的分子碰撞率依赖该速度，混合不同速度后的生存概率不一定是单一指数。上面的泊松模型用于平均率估算；把固定靶粒子的扫掠模型与全部粒子均运动的模型混用，会漏掉 $\sqrt2$。

\begin{example}[自由程与稀薄程度]\label{ex:freepath}
硬球气体的有效直径为 $d$，压强 $p$、温度 $T$。求平均自由程；在长为 $L$ 的管中，用平均碰撞模型估计分子穿越全管而不碰撞的概率。说明 $\lambda\gg L$ 与 $\lambda\ll L$ 各自意味着什么。
\Diagram{fig:freepathq}{管长与平均自由程的比较}{
 \draw[curve](0,0)rectangle(5,1);\draw[axis](.3,.5)--(4.6,.5);\draw[<->](0,-.4)--(5,-.4)node[midway,below]{$L$};
 \draw[guide](1,0)--(1,1);\draw[guide](2,0)--(2,1);\node at(1.5,1.4){典型长度 $\lambda$};
}
\end{example}
\begin{solution}
碰撞中心允许的最小距离为 $d$，截面 $\sigma=\pi d^2$。由 $n=p/(\kb T)$ 得
\[
 \lambda=\frac{\kb T}{\sqrt2\pi d^2p}.
\]
若用图\ref{fig:freepathq}中沿管轴的平均运动近似，取穿越时间 $t\simeq L/\bar v$，则 $P\simeq\ee^{-t/\tau}=\ee^{-L/\lambda}$。这不是对所有入射角、所有速度严格平均后的透射率；它衡量的是体碰撞是否重要。
定义 $\mathrm{Kn}=\lambda/L$，则这里的估计是 $P\simeq\ee^{-1/\mathrm{Kn}}$。当 $\mathrm{Kn}\gg1$，展开指数可得 $P\simeq1-L/\lambda$，一次体碰撞都不发生的粒子占绝大多数；当 $\mathrm{Kn}\ll1$，这类粒子的比例指数小。压强在定温下增大一倍，会使数密度和平均碰撞频率增大一倍，自由程减半。上述判断只涉及分子间碰撞；即使体碰撞很少，管壁的反射、吸附和孔道几何仍会影响传输。

\Diagram{fig:freepatha}{两种尺度下的典型轨迹}{
 \draw(0,0)rectangle(2.5,1);\draw[->](.15,.3)--(2.3,.7);\node at(1.25,-.4){$L/\lambda\ll1$};
 \draw(3.6,0)rectangle(6.1,1);\draw[->](3.75,.3)--(4.1,.7)--(4.4,.2)--(4.8,.8)--(5.1,.5)--(5.8,.7);\node at(4.85,-.4){$L/\lambda\gg1$};
}
图\ref{fig:freepatha}左图主要由壁面碰撞和几何边界决定输运；右图体碰撞多，局部平衡较容易建立。但“碰撞多”本身仍不能保证过程可逆，还需宏观变化足够慢、压强和温度梯度足够小。
\end{solution}

\section{活塞、分子流与壁面通量}
\begin{example}[自由管与活塞的相对速度]\label{ex:piston}
水平光滑面上有长管，管质量为 $m_2$，可无摩擦滑动的活塞质量为 $m_1$。管的右端开口，左端封闭，活塞最初位于中点。初始气体体积为 $V_0/2$、温度为 $T_0$、物质的量为 $\nu$。外界真空，装置起初静止。忽略气体质量相对于 $m_1,m_2$ 的贡献，气体可视为准静态绝热膨胀，$C_{V,m}$ 恒定。求活塞刚离开管口时相对于管的速度。
\Diagram{fig:pistonq}{初始半管气体与可反冲的管体}{
 \draw[curve](0,1.3)--(0,0)--(5,0);\draw[curve](0,1.3)--(5,1.3);\draw[line width=2pt](2.5,0)--(2.5,1.3);
 \node at(1.25,.65){$T_0,V_0/2$};\node[above]at(2.5,1.3){$m_1$};\node[below]at(1,-.15){管 $m_2$};\node at(4,.65){真空};
}
\end{example}
\begin{solution}
图\ref{fig:pistonq}的管体也可运动，不能把它当作固定边界。水平总动量为零；设两物体速度为 $v_1,v_2$，有 $m_1v_1+m_2v_2=0$。令 $u=v_1-v_2$，解得
\[
 v_1=\frac{m_2}{m_1+m_2}u,\qquad v_2=-\frac{m_1}{m_1+m_2}u.
\]
将两速度代入动能，得到
\[
 \frac12m_1v_1^2+\frac12m_2v_2^2
 =\frac12\frac{m_1m_2(m_1+m_2)}{(m_1+m_2)^2}u^2
 =\frac12\mu_pu^2,
 \quad \mu_p=\frac{m_1m_2}{m_1+m_2}.
\]
到管口时气体体积从半管增至全管。准静态绝热过程中，第一定律为
\[
 \nu C_{V,m}\dd T=-p\dd V
 =-\frac{\nu RT}{V}\dd V.
\]
除以$\nu C_{V,m}T$，把温度和体积分离后按实际端点积分：
\[
 \int_{T_0}^{T_f}\frac{\dd T}{T}
 =-\frac R{C_{V,m}}\int_{V_0/2}^{V_0}\frac{\dd V}{V},
 \quad [\ln T]_{T_0}^{T_f}
 =-\frac R{C_{V,m}}[\ln V]_{V_0/2}^{V_0}.
\]
于是$\ln(T_f/T_0)=-(R/C_{V,m})\ln2$，指数化得
\[
 T_f=T_0\,2^{1-\gamma},\qquad\gamma=1+\frac R{C_{V,m}},
 \quad T_0(V_0/2)^{\gamma-1}=T_fV_0^{\gamma-1}.
\]
这里的体积是管与活塞之间的相对几何体积，管体反冲不改变它由半管增至全管的事实。对整体使用能量守恒时，管与活塞所受气体力做的功都是内部能量交换，不再另加一次$\int p\dd V$。

\Diagram{fig:pistona}{离开管口前的能量和动量关系}{
 \draw[curve](0,0)rectangle(4.5,1);\draw[line width=2pt](4.5,0)--(4.5,1);\node at(2.1,.5){$T_f,V_0$};
 \draw[axis](4.5,1.4)--(5.7,1.4)node[right]{$v_1$};\draw[axis](1,1.4)--(-.2,1.4)node[left]{$v_2$};
 \node at(2.8,-.6){$\Delta U+\tfrac12\mu_pu^2=0$};
}
图\ref{fig:pistona}中内能减少转化为管和活塞的动能，故
\[
 \frac12\mu_pu^2=\nu C_{V,m}(T_0-T_f),\qquad
 u=\sqrt{\frac{2\nu C_{V,m}T_0}{\mu_p}(1-2^{1-\gamma})}.
\]
该结果以气体能及时重新平衡为前提。估计 $u^2/v_{\rm th}^2\sim\nu M/\mu_p$，所以 $\nu M\ll\mu_p$ 有利于活塞运动远慢于热运动；分子碰撞和声学松弛的时间也须远短于体积显著改变的时间。在这些条件下，气体才可沿一列平衡态演化并采用可逆绝热关系 $TV^{\gamma-1}=\text{常量}$。
\end{solution}

\subsection{运动壁面的碰撞积分}
\label{th:movingwall}
当 $L\ll\lambda$ 时，分子主要在壁面间自由飞行，瞬时速度分布通常带有历史信息。此时要计算碰撞次数及每次交换的动量，不能预先给气体套用平衡绝热线。先处理局部已知入射分布的壁面模型：平面壁沿 $+x$ 方向以速度 $U$ 运动，气体在壁左侧，入射分子必须满足 $v_x>U$。单位面积、单位时间入射数为 $n(v_x-U)f_x(v_x)\dd v_x$，弹性反射给予壁面的冲量为 $2m(v_x-U)$。因此
\begin{equation}\label{eq:movingwall}
 p_w(U)=2mn\int_U^\infty(v_x-U)^2f_x(v_x)\dd v_x.
\end{equation}
这是给定壁前入射分布时的瞬时力，不表示整个有限气体在后续时刻始终保持同一 Maxwell 分布。反射后的分子可能再次撞壁，须随边界更新分布。

\begin{example}[缓慢退让壁面的压强]\label{ex:movingwall}
气体入射分布在实验室系中保持温度 $T$、数密度 $n$ 的 Maxwell 分布。求运动壁面压强的闭式表达，并求 $|U|\ll\sqrt{\kb T/m}$ 时的一阶修正。
\Diagram{fig:movingwallq}{退让壁面只接收追得上它的分子}{
 \draw[line width=1.6pt](3,-.5)--(3,1.5);\node[particle]at(0,.4){};\draw[axis](0,.4)--(1.6,.4)node[midway,above]{$v_x$};
 \draw[axis](3,1.2)--(4.2,1.2)node[right]{$U$};\node at(1.4,-.5){$v_x>U$};
}
\end{example}
\begin{solution}
图\ref{fig:movingwallq}的碰撞速度是 $v_x-U$，它同时出现在通量和冲量中。设 $z=\sqrt\alpha U$，代入式\eqref{eq:movingwall}并令 $t=\sqrt\alpha v_x$，得
\[
 p_w=\frac{2mn}{\alpha\sqrt\pi}\int_z^\infty(t-z)^2\ee^{-t^2}\dd t.
\]
记$A_j=\int_z^\infty t^j\ee^{-t^2}\dd t$。零阶项按$\erfc$的定义为
$A_0=\sqrt\pi\erfc z/2$，一阶项直接使用指数原函数：
\[
 A_1=\lim_{L\to\infty}\left[-\frac12\ee^{-t^2}\right]_z^L
 =\frac12\ee^{-z^2}.
\]
二阶项分部积分，得到
\[
 \begin{aligned}
 A_2&=\left[-\frac t2\ee^{-t^2}\right]_z^\infty
 +\frac12\int_z^\infty\ee^{-t^2}\dd t\\
 &=\frac z2\ee^{-z^2}+\frac{\sqrt\pi}4\erfc z.
 \end{aligned}
\]
展开$(t-z)^2$，原积分为$A_2-2zA_1+z^2A_0$。代入三项并用$1/\alpha=2\kb T/m$，得到
\[
 p_w=n\kb T\left[(1+2z^2)\erfc z-\frac{2z}{\sqrt\pi}\ee^{-z^2}\right].
\]
分部积分的边界项在$z$处保留，在无穷远处为零。压强来自非负被积函数的积分，始终非负；一阶速度修正同时包含碰撞频率和单次冲量的变化。

\Diagram{fig:movingwalla}{壁面功率来自相对速度决定的力}{
 \draw[line width=1.5pt](2,0)--(2,1.5);\draw[axis](2,.4)--(3.4,.4)node[right]{$F=Ap_w$};
 \draw[axis](2,1.2)--(3.1,1.2)node[right]{$U$};\node at(.4,.7){气体};\node at(2.1,-.6){$\dot W=Ap_wU$};
}
在 $z=0$ 邻域，$\erfc z=1-2z/\sqrt\pi+O(z^3)$，所以
\[
 p_w=n\kb T\left[1-\frac4{\sqrt\pi}\sqrt{\frac m{2\kb T}}U
 +O\!\left(\frac{mU^2}{\kb T}\right)\right].
\]
退让 $U>0$ 时压强减小；迎面运动时增大。图\ref{fig:movingwalla}中壁所获功率为 $Ap_wU$，正负随运动方向改变。$U=0$ 时恢复 $n\kb T$，也是积分系数的直接校核。
\end{solution}

\subsection{壁面通量及余弦定律}
\label{th:wallflux}
取面元 $\dd S$，法线与速度夹角为 $\theta$。能够在 $\dd t$ 内撞上该面元的粒子，来自体积 $v\cos\theta\,\dd t\dd S$。各向同性分布的方向概率为 $\dd\Omega/(4\pi)$，且 $\dd\Omega=2\pi\sin\theta\dd\theta$。因此
\[
 \dd N=n g(v)\dd v\frac{\dd\Omega}{4\pi}
 v\cos\theta\dd t\dd S.
\]
只积入射半球 $0\le\theta\le\pi/2$，有
\[
 j=\frac{\dd N}{\dd S\dd t}=n\bar v\int_0^{\pi/2}\frac12\sin\theta\cos\theta\dd\theta
=\frac{n\bar v}{4}[\sin^2\theta]_0^{\pi/2}
 =\frac14n\bar v.
\]
单独给出方向时，撞壁样本的立体角密度为 $\cos\theta/\pi$；对方位角积分后的极角密度为 $2\sin\theta\cos\theta$。它在 $45^\circ$ 处最大，并不意味着单位立体角的通量在斜向最大。

\begin{example}[由撞壁冲量求气体压强]\label{ex:gaspressure}
用各向同性速率分布 $g(v)$ 计算理想弹性壁面的压强，再代入 Maxwell 分布，推出理想气体状态方程。
\Diagram{fig:gaspressureq}{入射与反射的法向动量变化}{
 \draw[line width=1.4pt](3,-1)--(3,1.6);\draw[->](.5,1.4)--(3,0);\draw[->](3,0)--(.5,-1.4);
 \draw[guide](1,0)--(4,0)node[right]{法线};\node at(1.8,.9){$\theta$};\node at(1.7,-1){$\theta$};
}
\end{example}
\begin{solution}
图\ref{fig:gaspressureq}中粒子法向动量从 $mv\cos\theta$ 变为 $-mv\cos\theta$，壁面得到 $2mv\cos\theta$ 的冲量。把该冲量乘入射数，并除以面积和时间，
\begin{align*}
 p&=n\int_0^\infty g(v)\dd v\int_0^{\pi/2}
 \frac12\sin\theta\dd\theta\,(v\cos\theta)(2mv\cos\theta)\\
 &=nm\langle v^2\rangle\int_0^{\pi/2}\sin\theta\cos^2\theta\dd\theta
 =\frac13nm\langle v^2\rangle.
\end{align*}
角积分中令$\mu=\cos\theta$，则$\dd\mu=-\sin\theta\dd\theta$，入射半球对应$\mu$从1降至0，所以
\[
 \int_0^{\pi/2}\sin\theta\cos^2\theta\dd\theta
 =-\int_1^0\mu^2\dd\mu
 =\left[\frac{\mu^3}3\right]_0^1=\frac13.
\]
径向积分留下的是二阶速率矩，而非平均速率的平方。对各向异性气体，沿$x$法向的壁面应使用$p_x=nm\langle v_x^2\rangle$；本题因三个方向等价，才有$\langle v_x^2\rangle=\langle v^2\rangle/3$，并可用一个标量压强描述各个壁面。

\Diagram{fig:gaspressurea}{通量中的法向投影与弹性反射冲量}{
 \draw[line width=1.3pt](2.5,-1.2)--(2.5,1.5);
 \draw[axis](.1,1.1)--(2.5,0)node[midway,above]{$\vct v$};\draw[axis](2.5,0)--(.1,-1.1);
 \draw[guide](.2,0)--(3.2,0);\draw[axis](.2,-1.5)--(2.2,-1.5)node[midway,below]{$v_n=v\cos\theta$};
 \draw(1.6,0)arc[start angle=180,end angle=155.4,radius=.9];\node at(1.36,.23){$\theta$};
 \node[right,align=left]at(3.6,.8){单位时间入射数\\含因子 $v_n$};
 \node[right,align=left]at(3.6,-.6){每次反射给壁面\\冲量 $2mv_n$};
}
图\ref{fig:gaspressurea}中的两个余弦分别来自扫掠体积和冲量；漏掉任何一个都会得到错误角因子。代入 $m\langle v^2\rangle=3\kb T$，得 $p=n\kb T$，再用 $N=\nu N_A$、$R=N_A\kb$，即 $pV=\nu RT$。
\end{solution}

\begin{example}[逸出粒子的速率与能量]\label{ex:effusivesample}
容器中气体处于温度 $T$ 的平衡态。壁上小孔尺寸远小于平均自由程，孔外真空，孔薄且不额外筛选方向。求穿孔粒子的归一化速率密度、平均速率和平均平动能。
\Diagram{fig:effusiveq}{小孔泄流与容器内的粒子抽样}{
 \draw[curve](0,0)--(0,2)--(3,2)--(3,1.15);\draw[curve](3,.85)--(3,0)--(0,0);
 \foreach \p in{(.4,.5),(.8,1.4),(1.6,.8),(2.3,1.5)}\node[particle]at\p{};
 \draw[axis](3,1)--(4.7,1)node[right]{真空};\node at(1.5,-.5){$T,n$};
}
\end{example}
\begin{solution}
图\ref{fig:effusiveq}中的小孔按“单位时间经过”抽样，同速率下快粒子有更多穿孔机会。角积分对各速率都给相同的 $1/4$，故归一化后的密度
\[
 g_e(v)=\frac{vg(v)}{\bar v}=2\alpha^2v^3\ee^{-\alpha v^2},\qquad
 \int_0^\infty g_e\dd v=2\alpha^2 I_3=1.
\]
于是
\[
 \bar v_e=\frac{\langle v^2\rangle}{\bar v}
 =\frac{3\pi}8\bar v,\qquad
 \langle v^2\rangle_e=2\alpha^2 I_5=\frac{4\kb T}{m}.
\]
加权后的密度仍按 $\dd v$ 归一化，平均值计算中的分母始终为同一个总通量。平均速率与平均能量分别涉及原分布的二阶、三阶矩，因而两者受快粒子的影响程度不同。逸出速率密度的峰值由 $3/v-2\alpha v=0$ 确定，对应 $v^2=3/(2\alpha)$；它比体内速率分布的峰向右移动。此处只改变了被观察的样本，容器在这一瞬时的温度仍是原来的 $T$，不能据束流平均能量另给它定义一个平衡温度。

\Diagram{fig:effusivea}{相同温度下体内分布与逸出分布的比较}{
 \draw[axis](0,0)--(5,0)node[right]{$v$};\draw[axis](0,0)--(0,2.3)node[above]{密度};
 \draw[curve]plot[domain=0:4.5,samples=70](\x,{4/sqrt(pi)*\x*\x*exp(-\x*\x)});
 \draw[curve,dashed]plot[domain=0:4.5,samples=70](\x,{2*\x*\x*\x*exp(-\x*\x)});
 \node at(3.3,1.9){实线 $g$};\node at(3.4,1.4){虚线 $g_e$（示意）};
}
逸出粒子的平均平动能为 $2\kb T$，高于体内的 $3\kb T/2$。图\ref{fig:effusivea}显示快粒子权重增加；容器若绝热且不补充气体，余气会冷却。冷却程度要进一步对变化的 $N,T$ 写能量守恒，不能只以“快粒子少了”代替计算。
\end{solution}

\subsection{相对变量独立性的温度条件}
\label{th:relativecorrelation}
质心与相对动能的代数分离对任何速度都成立，但它们在概率上独立，还依赖初始分布。考虑两个分别按温度 $T_1,T_2$ 独立制备的粒子群，在尚未通过碰撞建立共同温度的瞬间，任一方向均有
$\langle v_{1x}^2\rangle=\kb T_1/m_1$、
$\langle v_{2x}^2\rangle=\kb T_2/m_2$。
用第一讲已有的线性变换，
\begin{align*}
 \langle u_x^2\rangle&=\kb\left(\frac{T_1}{m_1}+\frac{T_2}{m_2}\right),\\
 \langle V_{cx}^2\rangle&=
 \frac{\kb(m_1T_1+m_2T_2)}{(m_1+m_2)^2},\\
 \langle V_{cx}u_x\rangle&=
 \frac{\kb(T_1-T_2)}{m_1+m_2}.
\end{align*}
最后一式由展开乘积后使用 $\langle v_{1x}v_{2x}\rangle=0$ 得到。不同温度时交叉协方差一般不为零，联合高斯密度不能再分解为质心、相对两个因子的乘积。等温时交叉项消失，对联合高斯变量，这又足以保证独立，恢复前文的分离结果。

\Diagram{fig:relativecorrelation}{同温与不同温度下联合分布等密度线的示意}{
 \begin{scope}
  \draw[axis](-1.45,0)--(1.6,0)node[right]{$V_{cx}$};
  \draw[axis](0,-1.1)--(0,1.3)node[above]{$u_x$};
  \draw[curve](0,0)ellipse(1 and .55);
  \node at(0,-1.6){$T_1=T_2$};
 \end{scope}
 \begin{scope}[xshift=5.5cm]
  \draw[axis](-1.45,0)--(1.6,0)node[right]{$V_{cx}$};
  \draw[axis](0,-1.1)--(0,1.3)node[above]{$u_x$};
  \draw[curve,rotate=30](0,0)ellipse(1.2 and .4);
  \node at(0,-1.6){$T_1>T_2$，存在正关联};
 \end{scope}
}
图\ref{fig:relativecorrelation}中的坐标各自按合适的速度尺度绘制；倾斜表示条件均值随另一变量改变。即使有这种关联，单独对质心变量积分后的相对速度边缘分布仍是高斯，其宽度就是上式的 $\langle u_x^2\rangle$。因此不同温度的相对速率习题可以只用边缘分布求平均碰撞率；若求碰撞后实验室系的能量分配，则还应保留质心与相对变量之间的关联。

在每个方向上，联合高斯变量的条件均值满足
$\langle V_{cx}\mid u_x\rangle
=\langle V_{cx}u_x\rangle u_x/\langle u_x^2\rangle$。
因此在非同温情况下，选择某个有符号的相对速度分量，也会改变样本的平均质心速度。这里 $T_1,T_2$ 描述制备或瞬时分布；持续碰撞会使两组分交换能量，最终建立共同温度。

\subsection{宏观漂移、热运动与动量通量}
\label{th:pressuretensor}
有整体流动时，实验室速度可分成平均流速与相对流体的热运动：
$\vct v=\vct V_0+\vct c$，其中 $\langle\vct c\rangle=0$。
局域平衡的 Maxwell 分布作用于 $\vct c$。温度由无规则运动的二阶矩确定，实验室系的总平均动能则为
\[
 \frac m2\langle v^2\rangle
 =\frac m2V_0^2+\frac m2\langle c^2\rangle.
\]
第一项是宏观流动能，不能计入温度。不同惯性系中的 $\vct V_0$ 改变，而 $\vct c$ 的统计宽度不变。

穿过法向为 $j$ 方向面元的粒子携带 $i$ 方向动量。每粒子动量为 $mv_i$，穿面通量含 $v_j$，故总动量通量张量为
\[
 \Pi_{ij}=nm\langle v_iv_j\rangle
 =\rho V_{0i}V_{0j}+P_{ij},\qquad
 P_{ij}=nm\langle c_ic_j\rangle,\quad \rho=nm.
\]
交叉项因 $\langle c_i\rangle=0$ 消失。各向同性平衡时 $P_{ij}=n\kb T\,\delta_{ij}$，才可用一个标量压强 $p$ 表示。若分布各向异性，三块不同朝向的壁可能承受不同的法向应力。

\begin{example}[各向异性粒子群的壁面压力]
\label{ex:anisotropicpressure}
长方体容器内有 $N$ 个质量为 $m$ 的独立粒子，速度分布关于各坐标反演对称，分量互不相关，且 $\langle v_i^2\rangle=s_i^2$。壁面均作镜面弹性反射。求三组壁面的压力及总平动能，并说明能否只用平均动能给任意一块壁确定压力。
\Diagram{fig:anisotropicpressureq}{三个方向的速度宽度与相应壁面}{
 \draw[curve](0,0)rectangle(3,2);
 \draw[curve](0,2)--(.9,2.8)--(3.9,2.8)--(3,2);
 \draw[curve](3,0)--(3.9,.8)--(3.9,2.8);
 \draw[axis](1.3,1)--(2.6,1)node[right]{$s_x$};
 \draw[axis](1.3,1)--(1.3,2.2)node[above]{$s_z$};
 \draw[axis](1.3,1)--(2,1.6)node[above left]{$s_y$};
 \node[right]at(4.4,1.3){$n=N/V$};
}
\end{example}
\begin{solution}
图\ref{fig:anisotropicpressureq}中，垂直 $x$ 轴的壁面只反转 $v_x$。对入射半空间求撞壁冲量，沿用前文的“通量乘单次冲量”，有
\[
 p_x=2nm\int_{v_x>0}v_x^2f_x(v_x)\dd v_x
 =nm s_x^2.
\]
第二步利用分量分布的反演对称性，把半轴二阶矩化成整轴的一半。相同方法给 $p_y=nm s_y^2$、$p_z=nm s_z^2$。总能量为
\[
 U_{\rm tr}=\frac{Nm}{2}(s_x^2+s_y^2+s_z^2)
 =\frac V2(p_x+p_y+p_z).
\]
\Diagram{fig:anisotropicpressurea}{法向二阶矩分别决定三组壁面压力}{
 \draw[axis](0,0)--(5,0)node[right]{壁面法向};
 \foreach \x/\h/\t in{1/1.1/x,2.5/1.8/y,4/.7/z}{
  \draw[fill=gray!15](\x-.3,0)rectangle(\x+.3,\h);
  \node[below]at(\x,0){$\t$};
  \node[above]at(\x,\h){$p_{\t}$};
 }
}
图\ref{fig:anisotropicpressurea}示意三方向不等的压力。给定 $U_{\rm tr}$ 只固定三个压力之和，仍可有不同的分配。只有各向同性时 $p_x=p_y=p_z=p$，才能进一步写成 $pV=2U_{\rm tr}/3$。若分子间碰撞逐渐使三方向能量均分，在体积固定且不向外散热的条件下，总能量不变，最终共同压力等于原先三个压力的算术平均。
\end{solution}

\subsection{碰撞事件的速率权重}
两组分在同温下的相对速率密度已由式\eqref{eq:relative}确定。某个相对速率 $u$ 对碰撞率的贡献，除概率密度外还乘扫掠速率 $u$ 和截面 $\sigma(u)$，故
\[
 K(T)=\langle\sigma u\rangle
 =\int_0^\infty\sigma(u)u\,g_u(u)\dd u.
\]
$K$ 的量纲为体积除时间，乘 $n_1n_2$ 后才成为单位体积、单位时间的异种碰撞数。按碰撞事件抽样的密度是 $\sigma(u)u g_u(u)/K$。截面随能量变化时，即使热平衡的 $g_u$ 未变，碰撞样本的分布也会改变。

\begin{example}[速率相关截面的碰撞平均]
\label{ex:powercrosssection}
两种粒子处于共同温度 $T$，约化质量为 $\mu$。在所讨论速率范围内，碰撞截面可近似为 $\sigma(u)=\sigma_\ast(u/u_\ast)^q$，其中 $u_\ast>0$、$q>-4$。求速率系数 $K(T)$，以及按碰撞事件抽样的平均相对动能。
\Diagram{fig:powercrosssectionq}{相对速率分布经扫掠与截面两次取权}{
 \node[draw](a)at(0,0){$g_u(u)\dd u$};
 \node[draw](b)at(3.1,0){$u\,\sigma(u)$};
 \draw[axis](a)--(b);
 \node at(1.55,.6){碰撞机会};
 \node[draw](c)at(6.2,0){事件样本};
 \draw[axis](b)--(c);
}
\end{example}
\begin{solution}
图\ref{fig:powercrosssectionq}中的权重为
$\sigma_\ast u^{q+1}/u_\ast^q$。用前文的一般 Maxwell 矩，把其中的质量换成 $\mu$，
\[
 K=\frac{\sigma_\ast}{u_\ast^q}
 \frac2{\sqrt\pi}\Gamma\!\left(\frac{q+4}{2}\right)
 \left(\frac{2\kb T}{\mu}\right)^{(q+1)/2}.
\]
积分原点附近的被积函数正比于 $u^{q+3}$，收敛要求 $q>-4$；无穷远端由高斯因子保证收敛。这个条件对应题设的指数范围。若真实截面只在有限速率区间服从幂律，则上式只适用于主要碰撞事件集中在该区间的温度。

事件平均必须用同一总碰撞率作分母，因此
\[
\begin{aligned}
 \langle E_r\rangle_{\rm coll}
 &=\frac{\mu}{2}\frac{\langle u^{q+3}\rangle}
 {\langle u^{q+1}\rangle}
 =\frac{\mu}{2}\frac{I_{q+5}(\alpha_u)}{I_{q+3}(\alpha_u)}\\
 &=\frac{\mu}{2}\frac{q+4}{2\alpha_u}
 =\frac{q+4}{2}\kb T,\qquad \alpha_u=\frac\mu{2\kb T}.
\end{aligned}
\]
第二行使用第\ref{th:integrals}节已由分部积分证明的递推式，取$s=q+3>-1$；两个速率矩具有相同归一化系数，作比后已消去。
\Diagram{fig:powercrosssectiona}{截面指数改变碰撞速率与事件能量}{
 \draw[axis](0,0)--(5,0)node[right]{$q$};
 \draw[axis](0,0)--(0,3)node[above]{$\langle E_r\rangle_{\rm coll}/(\kb T)$};
 \draw[curve](.6,.3)--(4.5,2.25);
 \fill(2,1)circle(1.2pt)node[below right]{$(-2,1)$};
 \fill(4,2)circle(1.2pt)node[above left]{$(0,2)$};
 \node[below]at(0,0){$-4$};
}
图\ref{fig:powercrosssectiona}中的直线只在 $q>-4$ 适用。常截面 $q=0$ 给出 $2\kb T$，比随机粒子对的 $3\kb T/2$ 大；$q=-1$ 时 $\sigma u$ 为常数，事件取样与原相对速率取样一致，平均能量恢复为 $3\kb T/2$。两个极限分别检验了碰撞加权和归一化。
\end{solution}

\subsection{局域平衡与边界过程的时间尺度}
一个容器中的总粒子数和总能量不足以确定随时间演化的分布。采用瞬时平衡的 Maxwell 形式，还要求碰撞与宏观变化之间有时间尺度分离。典型碰撞时间为 $\tau_c\sim\lambda/\bar v$，声学信息穿越尺度 $L$ 的时间约为 $\tau_s\sim L/c_s$。若体积或温度显著改变的时间为 $\tau_{\rm mac}$，准静态处理至少要求相关的松弛时间远小于 $\tau_{\rm mac}$。

小孔泄流还要求“孔附近的传输”与“容器内部的重新平衡”分别满足条件。孔径远小于平均自由程，可忽略孔口内的分子间碰撞；容器内却可通过足够多的体碰撞或壁面热化保持平衡。因而分子流小孔与平衡余气可以同时存在。若整个容器都处于无碰撞极限且壁面只作镜面反射，最初分布随边界运动会发生形变，须跟踪速度方向及粒子飞行历史。

宏观运动缓慢时，还应判断各处是否接近平衡。有限压差推动活塞、有限温差传热及流动耗散都会产生熵；降低平均速度未必能同时消除这些因素。后两讲将按具体能量交换方式和过程约束计算相应的熵产生。


\section{碰撞抽样与运动边界的进一步应用}
同一气体的体内样本、穿孔样本和碰撞样本具有不同的速率权重。体内样本按粒子数计数；穿孔和碰撞则按单位时间发生的事件计数，要在体内密度上乘相应的相对通量，再归一化。

当微观速度与宏观边界速度相差很大时，平衡压强可作为边界受力的首项；当二者可比时，要保留相对速度、入射阈值和反射规律。下面将这些条件分别展开。

\begin{example}[由球坐标变换核算方向概率]
三维速度分布各向同性，矢量密度为 $f(v)$。从球坐标的雅可比出发求 $(v,\theta,\phi)$ 的联合概率密度，证明速率与方向独立，并求 $v_z/v$ 的密度。
\Diagram{fig:angularjacq}{球坐标中的极角、方位角和速度分量}{
 \begin{scope}[x={(1cm,0)},y={(-.55cm,-.4cm)},z={(0,1cm)}]
 \draw[axis](0,0,0)--(2.5,0,0)node[right]{$v_x$};\draw[axis](0,0,0)--(0,2.3,0)node[below left]{$v_y$};\draw[axis](0,0,0)--(0,0,2.7)node[above]{$v_z$};
 \draw[axis](0,0,0)--(1.448,1.014,1.768)node[above right]{$\vct v$};
 \draw[guide](1.448,1.014,1.768)--(1.448,1.014,0);\draw[guide](0,0,0)--(1.448,1.014,0);
 \draw[curve]plot[domain=0:45]({.75*sin(\x)*cos(35)},{.75*sin(\x)*sin(35)},{.75*cos(\x)});
 \node at(.5,.15,.9){$\theta$};
 \draw[curve]plot[domain=0:35]({.8*cos(\x)},{.8*sin(\x)},0);\node at(.95,.35,0){$\phi$};
 \end{scope}
}
\end{example}
\begin{solution}
图\ref{fig:angularjacq}的坐标变换为 $v_x=v\sin\theta\cos\phi$、$v_y=v\sin\theta\sin\phi$、$v_z=v\cos\theta$。三个坐标微元方向彼此正交，其长度分别为 $\dd v$、$v\dd\theta$、$v\sin\theta\dd\phi$，故体积元为 $v^2\sin\theta\dd v\dd\theta\dd\phi$。于是
\[
 \dd P=f(v)v^2\dd v\sin\theta\dd\theta\dd\phi
 =g(v)\dd v\frac{\sin\theta\dd\theta\dd\phi}{4\pi}.
\]
右侧两个因子分别归一化，故速率与方向独立。方位角均匀分布；极角则按球带面积取权，密度为 $q(\theta)=\sin\theta/2$。
方向分布的归一化由$[\phi]_0^{2\pi}=2\pi$和$[-\cos\theta]_0^\pi=2$给出。令$\mu_z=\cos\theta=v_z/v$，则$|\dd\mu_z|=\sin\theta\dd\theta$，所以$\mu_z$在$[-1,1]$上有均匀密度$q_\mu=1/2$。其前两个矩为
\[
 \langle\mu_z\rangle=\frac12\left[\frac{\mu_z^2}{2}\right]_{-1}^1=0,
 \qquad \langle\mu_z^2\rangle=\frac12\left[\frac{\mu_z^3}{3}\right]_{-1}^1=\frac13.
\]
一阶矩为零说明没有优选方向；二阶矩$1/3$使速度平方的平均值均分到三个坐标方向。

\Diagram{fig:angularjaca}{相等极角宽度的球带面积不同}{
 \draw[curve](0,0)circle(1.6);\draw[curve](-1.6,0)--(1.6,0);\draw[guide](-1.38,.8)--(1.38,.8);\draw[guide](-.7,1.44)--(.7,1.44);\node[right]at(2,1){球带面积 $2\pi v^2\sin\theta\dd\theta$};
}
图\ref{fig:angularjaca}的球带面积随$\theta$变化。各向同性对应等立体角等概率；若把极角本身误设为均匀分布，便会得到错误的$\langle\cos^2\theta\rangle=1/2$。
\end{solution}

\begin{example}[固定速度粒子的相对碰撞频率]
背景气体分子速度大小均为 $w$，方向各向同性，数密度 $n$。一个试探粒子以固定速度大小 $v$ 运动，与背景粒子的硬球截面为常量 $\sigma$。求角度平均相对速率和试探粒子碰撞频率；分别讨论 $v>w$、$v<w$。
\Diagram{fig:fixedrelativeq}{固定试探速度与背景方向之间的夹角}{
 \draw[axis](0,0)--(3.2,0)node[right]{$\vct v$};\draw[axis](0,0)--(1.4,1.8)node[above]{$\vct w$};\draw[axis](1.4,1.8)--(3.2,0)node[midway,right]{$\vct v-\vct w$};\node at(.8,.35){$\theta$};
}
\end{example}
\begin{solution}
图\ref{fig:fixedrelativeq}中 $u=\sqrt{v^2+w^2-2vw\cos\theta}$。由于 $\cos\theta$ 在 $[-1,1]$ 上均匀，
\[
 \bar u(v,w)=\frac12\int_{-1}^1\sqrt{v^2+w^2-2vw\mu}\dd\mu.
\]
代换 $z=v^2+w^2-2vw\mu$，$\dd\mu=-\dd z/(2vw)$，积分端点为 $(v-w)^2$ 与 $(v+w)^2$，所以
\[
\begin{aligned}
\bar u(v,w)&=\frac1{4vw}\int_{(v-w)^2}^{(v+w)^2}\sqrt z\dd z
=\frac1{4vw}\left[\frac23z^{3/2}\right]_{(v-w)^2}^{(v+w)^2}\\
&=\frac{(v+w)^3-|v-w|^3}{6vw}.
\end{aligned}
\]
当$v>w>0$时，分子为$(v+w)^3-(v-w)^3=6v^2w+2w^3$；当$w>v>0$时，交换$v,w$即可。因此
\[\bar u(v,w)=\begin{cases}
v+w^2/(3v),&v\ge w>0,\\
w+v^2/(3w),&w\ge v>0.
\end{cases}\]
某一速率为零时取连续极限，分别为$\bar u(v,0)=v$和$\bar u(0,w)=w$；两者都为零时碰撞率为零。
同一变量变换也给出相对速率的允许区间：两速度平行时相对速率最小，反平行时最大，即 $|v-w|\le u\le v+w$。由 $2u\dd u=-2vw\dd\mu$，得到方向样本的密度 $q(u)=u/(2vw)$。积分区间的两端平方相差 $4vw$，因而该密度归一化。碰撞率还要再乘一个 $u$，所以即使背景速度大小固定，接近反向运动的粒子对也会在碰撞事件中占更高权重。

\Diagram{fig:fixedrelativea}{两种速度大小次序给出分段表达}{
 \node[draw,align=center]at(0,0){$v\ge w$\\$v+w^2/(3v)$};\node[draw,align=center]at(4.6,0){$w\ge v$\\$w+v^2/(3w)$};\node at(2.3,-1){$z(v)=n\sigma\bar u(v,w)$};
}
展开立方即得图\ref{fig:fixedrelativea}的分段式。在 $v=w$ 时两式均为 $4v/3$，连续衔接；$v\to0$ 时结果为 $w$，快速试探粒子 $v\gg w$ 时趋于 $v$。碰撞频率 $z(v)=n\sigma\bar u(v,w)$ 随试探粒子的速度而变。对 Maxwell 背景，还须用背景速率密度对 $w$ 再作一次平均。
\end{solution}

\begin{example}[按碰撞事件抽样的相对能量]
两个同温平衡组分相互独立，相对速度对应约化质量 $\mu$。碰撞截面为与速率无关的常量。求随机抽取一次碰撞时的相对速率密度、平均相对平动能，以及相对能量超过阈值 $E_a$ 的碰撞比例。
\Diagram{fig:collisionenergyq}{一次事件的权重含有相对扫掠速度}{
 \draw[curve](0,0)ellipse(.25 and .6);\draw[curve](0,.6)--(3,.6);\draw[curve](0,-.6)--(3,-.6);\draw[axis](.3,0)--(2.7,0)node[midway,above]{$u$};\node at(1.5,-1.1){事件率 $\propto u\sigma$};
}
\end{example}
\begin{solution}
相对速率的体样本密度为 $g_u(u)\propto u^2\ee^{-\alpha_u u^2}$，$\alpha_u=\mu/(2\kb T)$。图\ref{fig:collisionenergyq}中单位时间碰撞次数还乘 $u\sigma$，故事件样本密度为
\[
 g_{\rm coll}(u)=\frac{ug_u(u)}{\bar u}=2\alpha_u^2u^3\ee^{-\alpha_uu^2}.
\]
这与小孔泄流的速率加权结构相同，平均相对能量为 $\mu\langle u^2\rangle_{\rm coll}/2=2\kb T$。
把速率密度变为能量密度时，$\dd E=\mu u\dd u$。因此 $u^3\dd u=u^2\dd E/\mu$，速率三次因子在能量变量中变成一次因子。尾概率的积分可用
\[
 \int x\ee^{-x}\dd x=-(x+1)\ee^{-x}
\]
计算，下限贡献 $(x_a+1)\ee^{-x_a}$，上限为零。这里是碰撞发生以后再问其中多少超过阈值，分母已经是全部碰撞次数；若从瞬时任意粒子对中抽样，则应使用未经速率加权的相对能量分布。

\Diagram{fig:collisionenergya}{碰撞相对能量的归一化密度}{
 \draw[axis](0,0)--(5,0)node[right]{$E/(\kb T)$};\draw[axis](0,0)--(0,1.5)node[above]{$\kb T h_{\rm coll}$};\draw[curve]plot[domain=0:4.6](\x,{\x*exp(-\x)});\draw[guide](2.2,0)--(2.2,.8);\node[below]at(2.2,0){$E_a/(\kb T)$};
}
令 $E=\mu u^2/2$，变换密度得 $h_{\rm coll}(E)=E\ee^{-E/(\kb T)}/(\kb T)^2$。图\ref{fig:collisionenergya}阈值右侧积分，令 $x=E/(\kb T)$，
\[
 P(E>E_a)=\int_{x_a}^\infty x\ee^{-x}\dd x=(1+x_a)\ee^{-x_a},\quad x_a=E_a/(\kb T).
\]
所得比例对应“总相对动能超过阈值”的判据。若反应要求沿连心线的能量超过阈值，允许的碰撞参数范围也随能量变化，须进一步对反应截面取平均，下一例讨论这一情形。
\end{solution}

\begin{example}[带能量阈值的硬球反应截面]
两组分碰撞时，只有接触前沿两球连心线的相对动能超过 $E_a$ 才反应。硬球接触中心距为 $d$，相对速率为 $u$，碰撞参数为 $b$。求反应截面 $\sigma_r(u)$，再求同温 Maxwell 相对运动下的速率系数 $\langle\sigma_ru\rangle$。
\Diagram{fig:lineofcentersq}{接触点\(P\)处的法向分量与碰撞参数}{
 \draw[curve](3,0)circle(1);\draw[axis](0,.6)--(2.2,.6)node[midway,above]{$u$};
 \draw[guide](0,0)--(3,0);\draw[<->](.4,0)--(.4,.6)node[midway,left]{$b$};
 \draw[curve](3,0)--(2.2,.6)node[midway,below left]{$d$};\fill(2.2,.6)circle(1.1pt)node[above left]{$P$};
 \draw[guide](2.2,.6)--(3,.6);\draw[curve](2.75,.6)arc(0:-36.87:.55);\node at(2.94,.38){$\theta$};
}
\end{example}
\begin{solution}
图\ref{fig:lineofcentersq}接触时，入射方向与连心线夹角满足 $\sin\theta=b/d$，所以法向相对能量为 $E_n=\mu u^2\cos^2\theta/2=E(1-b^2/d^2)$。反应要求 $b^2<d^2(1-E_a/E)$，且先要有 $E\ge E_a$。因此
\[
 \sigma_r(E)=\pi d^2\left(1-\frac{E_a}E\right),\quad E\ge E_a;
 \qquad \sigma_r=0,\ E<E_a.
\]
入射中心的落点在截面平面上按面积计数，碰撞参数位于$b$到$b+\dd b$的环带面积为$2\pi b\dd b$。因此允许的截面为$\int_0^{b_{\max}}2\pi b\dd b=[\pi b^2]_0^{b_{\max}}=\pi b_{\max}^2$。阈值恰被达到时，只允许正碰这一零面积集合，截面从零连续开始；能量远大于阈值时，允许圆盘趋于整个硬球截面。下面将这个面积因子与各相对能量处的事件密度相乘。

\Diagram{fig:lineofcentersa}{允许反应的碰撞参数圆盘随能量增大}{
 \draw[curve](0,0)circle(1.4);\fill[gray!20](0,0)circle(.9);\draw[curve](0,0)circle(.9);\draw[axis](0,0)--(.9,0)node[midway,above]{$b_{\max}$};\node[right]at(1.7,0){全碰撞半径 $d$};
}
图\ref{fig:lineofcentersa}的有效截面已经包含碰撞几何筛选。将其乘 $u$ 后对体样本平均，也可利用上一例的事件能量密度：
\begin{align*}
 \langle\sigma_ru\rangle
 &=\pi d^2\bar u\int_{E_a}^\infty
 \left(1-\frac{E_a}E\right)\frac{E}{(\kb T)^2}\ee^{-E/(\kb T)}\dd E\\
 &=\pi d^2\bar u\int_{x_a}^\infty(x-x_a)\ee^{-x}\dd x\\
 &=\pi d^2\bar u[-(x+1-x_a)e^{-x}]_{x_a}^\infty
 =\pi d^2\bar u\ee^{-x_a}.
\end{align*}
这里指数因子与上一例的 $(1+x_a)\ee^{-x_a}$ 不同，差别来自允许反应的碰撞参数面积随能量而变化。反应体积率为 $n_1n_2\langle\sigma_ru\rangle$，同种反应的无序对计数还须除二。
\end{solution}

\begin{example}[撞壁样本的法向与切向能量]
壁面静止，气体处于温度 $T$ 的平衡。按撞上壁面的事件随机选一个分子，求其入射法向速度 $v_n>0$ 的密度，以及法向动能和两个切向动能的平均值。
\Diagram{fig:wallcomponentsq}{碰撞抽样只对法向速度加权}{
 \draw[line width=1.2pt](3,-1)--(3,1.8);\draw[axis](0,0)--(3,1.2)node[midway,above]{$\vct v$};\draw[axis](0,0)--(3,0)node[midway,below]{$v_n$};\draw[axis](3,0)--(3,1.2)node[right]{$v_t$};
}
\end{example}
\begin{solution}
图\ref{fig:wallcomponentsq}的壁面每秒截获厚度 $v_n\dd t$ 内的粒子，所以在原高斯因子上乘 $v_n$。归一化分母为 $\int_0^\infty v_nf_n(v_n)\dd v_n=1/(2\sqrt{\pi\alpha})$，得到
\[
 q_n(v_n)=2\alpha v_n\ee^{-\alpha v_n^2},\quad v_n\ge0.
\]
法向二阶矩为$\int_0^\infty v_n^2q_n(v_n)\dd v_n=2\alpha I_3=1/\alpha=2\kb T/m$，故平均法向动能为$\kb T$。联合事件密度仍分离为$q_n(v_n)f_y(v_y)f_z(v_z)$，归一化权重只含法向变量，两个切向高斯因子各自积分为1，平均动能仍各为$\kb T/2$。三方向相加为$2\kb T$；法向能量偏高来自事件筛选，未经过碰撞的体内分布并没有被壁面加热。

\Diagram{fig:wallcomponentsa}{逸出平均平动能按壁面方向分解}{
 \node[draw,align=center]at(0,0){法向\\$\kb T$};\node[draw,align=center]at(3,0){切向1\\$\kb T/2$};\node[draw,align=center]at(6,0){切向2\\$\kb T/2$};\node at(3,-1){总和 $2\kb T$};
}
图\ref{fig:wallcomponentsa}给出了逸出能量 $2\kb T$ 的另一路线。法向速度的平均值为 $2\alpha I_2=\sqrt\pi/(2\sqrt\alpha)$；它不同于从体内只筛选 $v_n>0$ 得到的条件平均，也不同于三维速率平均。
\end{solution}

\begin{example}[运动镜面的单次碰撞能量交换]
质量 $m$ 的粒子沿 $+x$ 方向以速度 $v$ 追上以速度 $U<v$ 退让的重壁，碰撞弹性，壁速度不变。求反射后的粒子速度、粒子能量减少与壁获得的功，并将结果与壁面受力积分相联系。
\Diagram{fig:singlewallq}{在实验室系与壁面系中比较碰撞}{
 \draw[line width=1.4pt](3,-.4)--(3,1.5);\draw[axis](0,.5)--(1.7,.5)node[midway,above]{$v$};\draw[axis](3,1.2)--(4.3,1.2)node[right]{$U$};\node at(1.5,-.6){壁面系入射速度 $v-U$};
}
\end{example}
\begin{solution}
先转到图\ref{fig:singlewallq}的壁面系，碰撞前相对速度 $v-U$，碰后变为 $-(v-U)$。变回实验室系，$v'=U-(v-U)=2U-v$。壁面获得的冲量为 $\Delta p_w=m(v-v')=2m(v-U)$，对应功 $\Delta W_w=U\Delta p_w$。
本题的壁速在碰撞前后保持不变，可由大质量壁面或外界驱动实现。若换成有限质量的自由活塞，就须联立两物体的动量、动能守恒方程，求碰后的粒子速度和活塞速度。

\Diagram{fig:singlewalla}{壁面系反射后变回实验室系}{
 \node[draw]at(0,0){$v-U$};\draw[axis](.8,0)--(2.2,0)node[midway,above]{反射};\node[draw]at(3.4,0){$-(v-U)$};\draw[axis](4.6,0)--(5.8,0)node[midway,above]{$+U$};\node[draw]at(6.5,0){$v'$};
}
直接算粒子能量减少，
\[
 \frac m2[v^2-(2U-v)^2]
 =\frac m2(4Uv-4U^2)=2mU(v-U)=\Delta W_w.
\]
该式保留$U^2$项，对满足追赶条件$v>U$的碰撞没有使用慢壁近似。$U>0$时分子失能，迎面运动的$U<0$情形则使分子增能。碰后$v'$可以仍为正：粒子在实验室系中继续向前，但速度小于壁速，因而与壁面分离。

将图\ref{fig:singlewalla}的单次能量交换乘事件率 $n(v-U)f_x(v)\dd v$，并在 $v>U$ 上积分，得到能量传递率 $Up_w(U)$。它与压强积分对应，反映了同一批碰撞中动量和能量的交换。
\end{solution}

\begin{example}[一维无碰撞气体的缓慢活塞膨胀]
粒子被限制在左端固定、右端缓慢退让的长为 $L(t)$ 的一维盒中。粒子间不碰撞，活塞速度 $U=\dot L$ 远小于粒子往返速度，左壁弹性反射。求单粒子速度大小随 $L$ 的变化，进而求均匀初始相位样本的轴向能量标度。与三维平衡气体绝热关系比较。
\Diagram{fig:ballisticpistonq}{粒子在两壁之间往返，仅运动壁改变能量}{
 \draw[line width=1.4pt](0,-.6)--(0,1);\draw[line width=1.4pt](4,-.6)--(4,1);\draw[<->](.3,.2)--(3.7,.2)node[midway,above]{$v$};\draw[axis](4,.8)--(5,.8)node[right]{$U$};\draw[<->](0,-1)--(4,-1)node[midway,below]{$L(t)$};
}
\end{example}
\begin{solution}
图\ref{fig:ballisticpistonq}中粒子每与退让活塞碰撞一次，速率由 $v$ 减为 $v-2U$，忽略 $U/v$ 的高阶修正。相邻两次撞活塞的间隔约为 $2L/v$，所以平均变化率
\[
 \frac{\dd v}{\dd t}\simeq\frac{-2U}{2L/v}=-\frac{\dot L}L v.
\]
设$L=L_0$时该粒子的速率为$v_0>0$，分离变量并积分，
\[\int_{v_0}^{v(L)}\frac{\dd s}{s}=-\int_{L_0}^{L}\frac{\dd\ell}{\ell},
\qquad[\ln s]_{v_0}^{v(L)}=-[\ln\ell]_{L_0}^{L}.\]
故$\ln[v(L)/v_0]=-\ln(L/L_0)$，即$v(L)=v_0L_0/L$。因此$vL$为缓慢变化下的不变量，单粒子轴向能量为$mv_0^2(L_0/L)^2/2$。
若初态为各向同性温度 $T_0$，而随后只有 $x$ 方向长度由 $L_0$ 改为 $L$，可把三方向二阶矩分别记为 $m\langle v_i^2\rangle=\kb T_i$。于是
\[
 T_x=T_0(L_0/L)^2,\qquad T_y=T_z=T_0,
\]
\[
 U_{\rm tr}=\frac{N\kb T_0}{2}\bigl[(L_0/L)^2+2\bigr].
\]
这些 $T_i$ 只是分方向能量参数。它们不同，表明分布已偏离单温度 Maxwell 平衡态；横向能量不会自动补给不断损失轴向能量的分子。

\Diagram{fig:ballisticpistona}{缓慢循环碰撞给出轴向能量的平方反比律}{
 \draw[axis](0,0)--(5,0)node[right]{$L/L_0$};\draw[axis](0,0)--(0,2.8)node[above]{$E_x/E_{x0}$};\draw[curve]plot[domain=1:4.5](\x,{2.4/\x^2});\node[below]at(1,0){$1$};\node[left]at(0,2.4){$1$};
}
图\ref{fig:ballisticpistona}的标度也适用于对初始相位平均的独立粒子样本，但要求变化慢于各粒子的往返时间。非常慢的尾部粒子可能违反这一条件，不能把近似无限外推。

三维无碰撞盒只沿 $x$ 膨胀时，$E_x\propto L^{-2}$，横向 $E_y,E_z$ 不变，分布因而具有各向异性，需要分别描述三个方向的能量尺度。若碰撞能及时重新分配能量，气体才恢复单一平衡温度，并满足单原子气体的 $TV^{2/3}=\text{常量}$。边界运动缓慢与内部达到热平衡须分别判断。
\end{solution}

\begin{example}[时间飞行测温中的两种探测权重]
一个小孔发出的短脉冲分子束沿法线传播，速度分布取逸出速率密度 $g_e(v)$。探测器在距离 $L$ 处，接受极窄角度，对到达的每个粒子等效计数，忽略碰撞。求到达时间密度及峰值；若探测器测量固定时刻某薄层中的粒子数，权重有何变化？另比较同一速率谱的稳态束流在薄层中的瞬时计数。
\Diagram{fig:tofthermq}{短脉冲穿过固定飞行距离后到达计数器}{
 \draw[line width=1.2pt](0,-.6)--(0,1);\draw[line width=1.2pt](5,-.6)--(5,1);\draw[axis](.2,.2)--(4.7,.2)node[midway,above]{$v=L/t$};\draw[<->](0,-1)--(5,-1)node[midway,below]{$L$};\node at(0,1.4){孔};\node at(5,1.4){计数器};
}
\end{example}
\begin{solution}
图\ref{fig:tofthermq}中速度与时间一一对应，$v=L/t$、$|\dd v/\dd t|=L/t^2$。窄角度筛选不改变速率分布，因此到达时间密度
\[
 q_t(t)=g_e(L/t)\frac L{t^2}
 =2\alpha^2L^4t^{-5}\exp(-\alpha L^2/t^2),\quad t>0.
\]
对数求导，$\dd\ln q_t/\dd t=-5/t+2\alpha L^2/t^3$，峰值满足 $t_p^2=2\alpha L^2/5$，所以 $t_p=L\sqrt{m/(5\kb T)}$。到达时间密度的峰值还受雅可比 $L/t^2$ 影响，因此应直接对 $q_t$ 求极值，而不能把速率峰值简单取倒数。
到达时间密度的归一化可直接用 $v=L/t$ 验证：当 $t$ 从0增至无穷大，$v$ 从无穷大降至0，雅可比的绝对值抵消积分方向的反转，故总概率仍为1。长时间尾部有 $q_t\propto t^{-5}$，来自束流中很慢的分子；极短时间则被指数因子强烈抑制。仪器若按粒子到达事件计数，每个粒子只记一次；这与在固定薄层内持续观察、因停留时间而多看见慢粒子的情形不同。

\Diagram{fig:toftherma}{时间分布的峰值由雅可比与指数共同决定}{
 \draw[axis](0,0)--(5.4,0)node[right]{$t$};\draw[axis](0,0)--(0,2.8)node[above]{$q_t$};\draw[curve]plot[domain=.25:4.8,samples=90](\x,{2*exp(-1/(\x*\x))/\x^5});\node at(3.7,2.4){事件计数};
}
图\ref{fig:toftherma}是脉冲的到达时间密度。固定时刻的空间计数须另用$x=vt$作变量变换。把短脉冲的发射历时近似为零，接受方向内共含$N_p$个粒子；固定$t>0$时，$\dd x=t\dd v$，因此
\[
 \dd N=N_pg_e(v)\dd v
 =\frac{N_p}{t}g_e(x/t)\dd x.
\]
对图\ref{fig:toftherm-layer-a}中以$L$为中心、厚度为$\Delta\ell$的薄层，瞬时粒子数为
\[
 N_{\rm layer}(t)
 =N_p\int_{(L-\Delta\ell/2)/t}^{(L+\Delta\ell/2)/t}g_e(v)\dd v
 \simeq\frac{N_p\Delta\ell}{t}g_e(L/t).
\]
末式要求该窄速度区间内$g_e$变化很小。固定$t$和$L$已把速率限制在$L/t$附近，不能把整个脉冲的速率密度直接乘$1/v$来表示这一快照。

\Diagram{fig:toftherm-layer-a}{脉冲快照按位置筛选；稳态薄层计数另含停留时间}{
 \node[anchor=west]at(0,2.9){脉冲：固定$t$};
 \draw[axis](0,1.7)--(6.3,1.7)node[right]{$x$};
 \fill[black!8](3.4,1.8)rectangle(4.4,2.5);
 \draw[dashed](3.4,1.6)--(3.4,2.55);\draw[dashed](4.4,1.6)--(4.4,2.55);
 \foreach \x in{1.1,2.2,3.9,5.2}\fill(\x,2.15)circle(1.5pt);
 \draw[<->](3.4,2.7)--(4.4,2.7)node[midway,above]{$\Delta\ell$};
 \node[below]at(3.9,1.7){$L$};\node[above]at(1.6,2.15){$x=vt$};
 \node[anchor=west]at(0,.65){稳态：连续发射};
 \draw[axis](0,-.15)--(6.3,-.15)node[right]{$x$};
 \fill[black!8](3.4,-.05)rectangle(4.4,.45);
 \draw[axis](2.3,.2)--(4.9,.2)node[right]{$v$};
 \node[below]at(3.9,-.4){停留时间$\Delta\ell/v$};
}
稳态束流连续补充各速率的粒子。若总发射率为$J$，速率区间内的到达率为$Jg_e(v)\dd v$，每个粒子在薄层中停留$\Delta\ell/v$，故
\[
 \dd N_{\rm steady}=Jg_e(v)\dd v\frac{\Delta\ell}{v},\qquad
 q_{\rm space}(v)=\frac{g_e(v)/v}{\int_0^\infty g_e(w)\dd w/w}=g(v).
\]
最后用了$g_e(v)=vg(v)/\bar v$及$\int_0^\infty g(v)\dd v=1$。这里的$1/v$来自稳态停留时间；若对一个脉冲在薄层中的占据数作全时间累计，也会得到相同的停留时间权重。实际测温还须按装置条件计入角接受范围、孔道传输和仪器响应。
\end{solution}

这些计算中的权重，都能追溯到图中的扫掠体积或变量变换。先按面元和相对速度写出事件数，再对实际允许的速度区域积分，便能确定相应的平均量。

\section{自由飞行、方向随机化与扩散尺度}
\subsection{一次碰撞的抽样怎样变成长期输运}
先用恒速随机飞行模型求扩散系数。粒子速率为$v$，碰撞时刻构成平均间隔$\tau$的泊松过程，每次碰撞后方向独立且各向同性，平均自由程为$\ell=v\tau$。这一模型保留了自由飞行和方向随机化，但省去了Maxwell气体碰撞中速率的变化。

初始方向取各向同性。经过时间$t$没有碰撞的概率为$\ee^{-t/\tau}$；只要发生过至少一次方向随机化，当前速度与初始速度的平均点积便为零。因而速度相关函数为
\[
 C_v(t)=\langle\vct v(0)\cdot\vct v(t)\rangle=v^2\ee^{-t/\tau}.
\]
位移为$\vct r(t)=\int_0^t\vct v(s)\dd s$，将平方展开并取平均，得到$\langle r^2(t)\rangle=\int_0^t\int_0^t C_v(|s_1-s_2|)\dd s_1\dd s_2$。积分正方形中两个三角形的贡献相等；在$s_1\ge s_2$的三角形内，固定时间差$s=s_1-s_2$后，$s_2$从$0$到$t-s$，因此
\begin{align*}
 \langle r^2(t)\rangle
 &=2\int_0^t(t-s)C_v(s)\dd s\\
 &=2v^2\int_0^t(t-s)\ee^{-s/\tau}\dd s\\
 &=2v^2\left[\tau(s+\tau-t)\ee^{-s/\tau}\right]_0^t\\
 &=2v^2[\tau t-\tau^2(1-\ee^{-t/\tau})].
\end{align*}
短时间$t\ll\tau$时展开指数，首个非零项为$v^2t^2$，对应弹道运动；长时间$t\gg\tau$时为$2v^2\tau t=6Dt$，所以
\[
 D=\frac{v^2\tau}{3}=\frac{v\ell}{3}.
\]
系数$1/3$来自三维各向同性，不应移到一维仅左右随机的模型中照用。若速率变化、散射保留方向记忆或不同速率具有不同碰撞时间，必须重新求相关函数，不能把任何平均速率与平均自由程任意相乘。
\ExpansionDiagram{fig:dep-th1-walk}{短时间近似直线飞行，长期由许多随机方向段累积}{
\draw[->,thick](0,0)--(1,.4)--(1.6,1.3)--(2.6,1)--(3,.2)--(4,.7)--(4.4,1.8)--(5.2,1.4);\fill(0,0)circle(1.5pt)node[below]{$O$};\draw[dashed,->](0,0)--(5.2,1.4)node[midway,below]{$\vct r$};\node[above]at(2.5,2.2){每次碰撞后方向重新抽样};}
图\ref{fig:dep-th1-walk}中的实际路程始终为$vt$，端点位移却因方向抵消而只有典型大小$\sqrt{6Dt}$。路程与位移不能混用来估计扩散所需时间。

\subsection{从粗粒化位移到扩散方程}
在比$\tau$长、但比宏观扩散时间短的区间$\Delta t$中，单个方向满足$\langle\Delta x\rangle=0$、$\langle\Delta x^2\rangle=2D\Delta t$。设位移$\xi$的归一化分布为$P(\xi)$，且介质均匀、没有平均漂移，则
\[
 n(x,t+\Delta t)=\int_{-\infty}^{\infty}P(\xi)n(x-\xi,t)\dd\xi.
\]
在密度变化尺度上展开两边，保留时间一阶项与位移二阶项，得
\begin{align*}
 n+\Delta t\,\partial_tn
 &=n-\langle\xi\rangle\partial_xn
   +\frac12\langle\xi^2\rangle\partial_x^2n\\
 &=n+D\Delta t\,\partial_x^2n.
\end{align*}
约去$n$并除以$\Delta t$，再与粒子数守恒式$\partial_tn=-\partial_xj_N$比较；取均匀密度时通量为零，得到
\[
 \partial_t n=D\partial_x^2n,\qquad j_N=-D\partial_xn.
\]
其中$j_N$是粒子数通量。该连续近似要求密度变化尺度远大于$\ell$；壁面附近一个自由程内常有非各向同性层，边界条件应与吸收、反射或穿透的具体模型相符。

\begin{example}[两吸收端之间的平均逃逸时间]
粒子在$0<x<L$中扩散，两端到达即吸收。扩散系数为$D$，初始位置为$x_0$且远离壁面的微观边界层。求先从右端逃逸的概率，以及到达任一端的平均时间。
\ExpansionDiagram{fig:dep-th1-exit-q}{两吸收端的首次到达问题，初始点为$x_0$}{
\draw[->](0,0)--(6,0)node[right]{$x$};\draw[very thick](.5,-.6)--(.5,.8)node[above]{吸收};\draw[very thick](5.5,-.6)--(5.5,.8)node[above]{吸收};\fill(2.3,0)circle(2pt)node[below]{$x_0$};\node[below]at(.5,-.7){0};\node[below]at(5.5,-.7){$L$};\draw[<->](1.5,.4)--(3.1,.4);}
\end{example}
\begin{solution}
用$p(x)$表示右端先到达的概率。图\ref{fig:dep-th1-exit-q}中从内点经过短粗粒化时间后，后续概率等于$p(x+\Delta x)$的平均。于是
\[
 p(x)=p(x)+p'(x)\langle\Delta x\rangle
       +\frac12p''(x)\langle\Delta x^2\rangle
     =p(x)+D\Delta t\,p''(x),
\]
故$p''=0$。逐次积分得$p'=A$、$p=Ax+B$。从左端出发已经被左端吸收，所以$p(0)=0$，给出$B=0$；从右端出发已经到达右端，所以$p(L)=1$，给出$A=1/L$。因此$p(x_0)=x_0/L$。

用$T(x)$表示剩余平均到达时间。时间已经经过$\Delta t$，所以$T(x)=\Delta t+\langle T(x+\Delta x)\rangle$。展开右边得$T=\Delta t+T+D\Delta t\,T''$，约去$T$并除以$\Delta t$，得到$DT''=-1$。逐次积分为
\[
 T'(x)=-\frac{x}{D}+C_1,\qquad
 T(x)=-\frac{x^2}{2D}+C_1x+C_2.
\]
两端一旦到达就不再等待，故$T(0)=T(L)=0$。前一条件给出$C_2=0$，后一条件给出$C_1=L/(2D)$，因此
\[
 T(x)=\frac{x(L-x)}{2D}.
\]
\ExpansionDiagram{fig:dep-th1-exit-a}{右端到达概率为直线，平均等待时间为抛物线}{
\draw[->](0,0)--(5,0)node[right]{$x/L$};\draw[->](0,0)--(0,2.7)node[above]{归一化量};\draw[dotted](0,2)--(4,2);\draw(-.06,2)--(.06,2)node[left,xshift=-2pt]{$1$};\draw[dashed](0,0)--(4,2);\draw[thick,domain=0:4,samples=70]plot(\x,{2*\x*(4-\x)/4});\node[above]at(2,2){$T/T_{\max}$};\node[right]at(4,2){$p$};\node[below]at(4,0){1};}
图\ref{fig:dep-th1-exit-a}中最长平均等待发生在中央，$T_{\max}=L^2/(8D)$，而非$L/v$。在本节的方向随机化模型下，这等于$3L^2/(8v\ell)$；$L\gg\ell$时远大于一次直线穿越时间。若初态距壁面仅约一个自由程且方向已经给定，最初一段弹道运动不可由此扩散边界模型准确描述。
\end{solution}

\begin{example}[封闭容器中的浓度均匀化]
一长直容器截面积$A$、长度$L$，两端完全反射，初始密度$n(x,0)=n_0[1+\epsilon\cos(\pi x/L)]$，$|\epsilon|<1$。取恒定$D$，求密度、粒子流及不均匀振幅降为原来的$1/10$所需时间。
\ExpansionDiagram{fig:dep-th1-mix-q}{反射端无粒子流，初态浓度由左向右平滑变化}{
\draw(0,0)rectangle(5,1.6);\draw[very thick](0,0)--(0,1.6);\draw[very thick](5,0)--(5,1.6);\foreach\x/\n in{.5/5,1.5/4,2.5/3,3.5/2,4.5/1}{\foreach\j in{1,...,\n}{\fill(\x,{.2+1.2*\j/(\n+1)})circle(1.2pt);}}\node[below]at(2.5,0){$\partial_xn=0$于两端};}
\end{example}
\begin{solution}
图\ref{fig:dep-th1-mix-q}的无流边界是$\partial_xn=0$。设$n=n_0[1+a(t)\cos(\pi x/L)]$，初值给出$a(0)=\epsilon$。对空间变量逐次求导，得
\[
 \partial_xn=-\frac{\pi n_0a}{L}\sin\frac{\pi x}{L},\qquad
 \partial_x^2n=-\frac{\pi^2n_0a}{L^2}\cos\frac{\pi x}{L}.
\]
第一式在$x=0,L$都为零，满足反射边界。代入扩散方程并比较余弦项的系数，得到$\dot a=-D\pi^2a/L^2$。若$\epsilon\ne0$，分离变量并代入初值，
\[
 \int_\epsilon^{a(t)}\frac{\dd a}{a}
 =-\frac{D\pi^2}{L^2}\int_0^t\dd s,
 \qquad
 \ln\frac{|a(t)|}{|\epsilon|}=-\frac{D\pi^2t}{L^2}.
\]
振幅保持初始符号，故$a(t)=\epsilon\ee^{-\pi^2Dt/L^2}$；$\epsilon=0$时$a(t)=0$也满足方程与初值。再由$j_N=-D\partial_xn$得到
\[
 n(x,t)=n_0\left[1+\epsilon\ee^{-\pi^2Dt/L^2}\cos\frac{\pi x}L\right],
 \quad j_N=\frac{\pi Dn_0\epsilon}{L}\ee^{-\pi^2Dt/L^2}\sin\frac{\pi x}L.
\]
\ExpansionDiagram{fig:dep-th1-mix-a}{常数模保持总粒子数，余弦不均匀模随时间衰减}{
\draw[->](0,0)--(5.5,0)node[right]{$x/L$};\draw[thick,domain=0:5,samples=80]plot(\x,{1.3+.7*cos(180*\x/5)});\draw[dashed,domain=0:5,samples=80]plot(\x,{1.3+.25*cos(180*\x/5)});\draw[dotted](0,1.3)--(5,1.3);\node[right]at(5,1.3){$n_0$};}
图\ref{fig:dep-th1-mix-a}中常数模保持不变。将密度对容积积分，
\[
 N(t)=n_0A\left[x+\frac{\epsilon L}{\pi}
 \ee^{-\pi^2Dt/L^2}\sin\frac{\pi x}{L}\right]_0^L=n_0AL,
\]
可见余弦项不改变总粒子数。当$\epsilon\ne0$时，振幅降为原来的$1/10$要求$\ee^{-\pi^2Dt/L^2}=1/10$，取对数得$-\pi^2Dt/L^2=-\ln10$，所以$t=L^2\ln10/(\pi^2D)$。若$\epsilon=0$，初态已经均匀。若只使容器长度变为两倍而材料和温度不变，均匀化时间增为四倍。反射边界保留常数模，吸收边界则不保留粒子总数；两类问题不能共用同一组末态条件。
\end{solution}

\section{习题与解答}
\begin{exercise}[两种温度下的曲线缩放]
同种气体由 $T$ 变为 $4T$，粒子数不变。求速率分布峰值位置、峰高和区间概率 $P(a<v<b)$ 的变换关系。
\Diagram{fig:scaleq}{不同温度下分布变宽而面积保持为一}{
 \draw[axis](0,0)--(5.6,0)node[right]{$v$};\draw[axis](0,0)--(0,2.4)node[above]{$g$};
 \draw[curve]plot[domain=0:5,samples=80](\x,{4/sqrt(pi)*\x*\x*exp(-\x*\x)});
 \draw[curve,dashed]plot[domain=0:5,samples=80](\x,{1/(2*sqrt(pi))*\x*\x*exp(-\x*\x/4)});\node at(4,1.9){$T\to4T$};
}
\end{exercise}
\begin{solution}
设原密度为 $g_T(v)$，直接代入式\eqref{eq:speeddensity}，有 $g_{4T}(v)=\tfrac12g_T(v/2)$。图\ref{fig:scaleq}中横向伸长两倍、纵向降低一半，面积因而不变。峰值位置 $v_p\to2v_p$，峰高减半。
在变换后的密度中，温度的 $T^{-3/2}$ 因子给出 $1/8$，指数中的 $v^2/T$ 则通过 $v/2$ 保持原形；同时原密度内的速度平方又产生因子4，合起来才是前面的 $1/2$。区间积分另有 $\dd v=2\dd w$，正好与纵向缩放抵消。图中的两段区间代表同一组无量纲端点 $v/\sqrt T$；固定实验室速率端点时，端点相对于热运动尺度的位置已变化，因此其概率不会普遍保持不变。

\Diagram{fig:scalea}{概率区间随无量纲速度一起变换}{
 \draw[axis](0,1.2)--(5,1.2)node[right]{$T$};\draw[axis](0,0)--(5,0)node[right]{$4T$};
 \draw[line width=2pt](1,1.2)--(2,1.2);\draw[line width=2pt](2,0)--(4,0);
 \node[above]at(1,1.2){$a$};\node[above]at(2,1.2){$b$};\node[below]at(2,0){$2a$};\node[below]at(4,0){$2b$};
}
令 $w=v/2$，则 $P_{4T}(a<v<b)=\int_{a/2}^{b/2}g_T(w)\dd w$。图\ref{fig:scalea}表示的是 $P_{4T}(2a<v<2b)=P_T(a<v<b)$；若区间端点不跟随温度缩放，概率通常会改变。
\end{solution}

\begin{exercise}[阈值以上的速率]
求平衡气体中 $v>v_0$ 的比例，并令 $v_0=v_p$ 写成误差函数形式。
\Diagram{fig:speedtailq}{速率阈值右侧的概率}{
 \draw[axis](0,0)--(5,0)node[right]{$v$};\draw[axis](0,0)--(0,2)node[above]{$g$};
 \draw[curve]plot[domain=0:4.5,samples=70](\x,{3*\x*\x*exp(-\x*\x)});\draw[guide](1.5,0)--(1.5,1.7);\node[below]at(1.5,0){$v_0$};
}
\end{exercise}
\begin{solution}
令$t=\sqrt\alpha v$，则$\dd v=\dd t/\sqrt\alpha$，下限$v_0$变为$z=\sqrt\alpha v_0$、上限仍为无穷大。图\ref{fig:speedtailq}右侧概率因而为$(4/\sqrt\pi)\int_z^\infty t^2\ee^{-t^2}\dd t$。分部积分取$t$与$t\ee^{-t^2}\dd t$，得到
\[
 P(v>v_0)=\erfc z+\frac{2z}{\sqrt\pi}\ee^{-z^2}.
\]
分部积分的具体边界项为 $[-t\ee^{-t^2}/2]_z^\infty=z\ee^{-z^2}/2$，剩余积分为原高斯尾的一半。乘上速率密度归一化常数后，便分别形成图中的两项。阈值取最概然速率时 $\alpha v_p^2=1$，所以应把无量纲下限取为1。最概然值只给出密度峰的位置，不把总体分成各占一半的两部分；本题的积分一般也不等于 $1/2$。

\Diagram{fig:speedtaila}{速率尾积分由边界项和高斯尾组成}{
 \node[draw]at(0,0){边界项 $2z\ee^{-z^2}/\sqrt\pi$};\node at(3,0){$+$};\node[draw]at(5,0){$\erfc z$};
}
当 $v_0=v_p$ 时 $z=1$，结果为 $\erfc1+2/(\ee\sqrt\pi)$。图\ref{fig:speedtaila}中的边界项不能省去；速率密度有 $v^2$ 因子，与一维速度分量的尾概率不同。$z=0$ 给出1，$z\to\infty$ 给出0。
\end{solution}

\begin{exercise}[不同温度的相对速度]
两个独立粒子质量为 $m_1,m_2$，各自服从温度 $T_1,T_2$ 的无漂移 Maxwell 分布。求相对速率的平均值和均方值。能否仍把 $T_1$ 或 $T_2$ 直接代入约化质量公式？
\Diagram{fig:twotq}{两个独立速度分布的差}{
 \node[draw,align=center]at(0,0){$m_1,T_1$\\$\vct v_1$};\node[draw,align=center]at(3.4,0){$m_2,T_2$\\$\vct v_2$};\draw[->](.8,-.9)--(2.6,-.9)node[midway,below]{$\vct u=\vct v_1-\vct v_2$};
}
\end{exercise}
\begin{solution}
每个方向的速度分量均为零均值高斯，方差分别为 $\kb T_1/m_1$、$\kb T_2/m_2$。图\ref{fig:twotq}中的两者独立，差的方差相加，记 $s^2=\kb T_1/m_1+\kb T_2/m_2$，有 $u_x,u_y,u_z$ 独立且方差均为 $s^2$。
方差相加之前先展开 $u_x^2=v_{1x}^2+v_{2x}^2-2v_{1x}v_{2x}$。独立且均值均为零时，交叉项平均为 $\langle v_{1x}\rangle\langle v_{2x}\rangle=0$。三个方向具有相同方差，所以均方相对速率是三个分量二阶矩之和。若两组分还存在不同的宏观漂移速度，交叉项的处理要先减去各自均值，相对分布的中心也会偏移；本题给出的零漂移条件正是中心位于原点的依据。

\Diagram{fig:twota}{相对速度各分量的方差相加}{
 \draw[axis](-2.7,0)--(2.7,0)node[right]{$u_x$};\draw[curve]plot[domain=-2.5:2.5,samples=60](\x,{1.7*exp(-\x*\x/2)});\draw[<->](-1,-.4)--(1,-.4)node[midway,below]{$2s$};
}
故三维相对速率仍为 Maxwell 形状，$g_u(u)=\sqrt{2/\pi}\,u^2s^{-3}\ee^{-u^2/(2s^2)}$，于是 $\bar u=\sqrt{8/\pi}\,s$、$\langle u^2\rangle=3s^2$。图\ref{fig:twota}的宽度由两个温度共同决定，计算时应保留 $\kb T_1/m_1+\kb T_2/m_2$ 的完整组合；当 $T_1=T_2=T$ 时，才简化为 $s^2=\kb T/\mu$。
\end{solution}

\begin{exercise}[同种气体的碰撞计数]
体积 $V$ 内有 $N$ 个相同硬球，温度 $T$、有效截面 $\sigma$。求总体每秒碰撞次数，并比较用“每粒子的频率乘粒子数”得到的结果。
\Diagram{fig:paircountq}{同一次碰撞涉及两个粒子}{
 \node[circle,draw,minimum size=9mm](a)at(0,0){$i$};\node[circle,draw,minimum size=9mm](b)at(2,0){$j$};\draw[->](a)--(.7,0);\draw[->](b)--(1.3,0);\node at(1,-.8){一对 $\{i,j\}$};
}
\end{exercise}
\begin{solution}
单个粒子的平均碰撞频率为 $z=n\sigma\sqrt2\bar v$。图\ref{fig:paircountq}中一次碰撞同时出现在 $i$ 和 $j$ 的计数中，因此总体碰撞次数率为
\[
\mathcal R=\frac12Nz\simeq\frac{N^2\sigma}{\sqrt2V}\sqrt{\frac{8\kb T}{\pi m}}.
\]
单粒子频率$z$按参与者计数，$Nz$便把每次碰撞计入两次。图中合并这两条有序记录时，截面和相对运动的平均不变；有限粒子数还须排除与自身配对的项。

\Diagram{fig:paircounta}{有序计数转换为无序粒子对}{
 \node[draw]at(0,0){$i\to j$};\node[draw]at(2,0){$j\to i$};\draw[->](0,-.4)--(1,-1);\draw[->](2,-.4)--(1,-1);\node[below]at(1,-1){同一事件，除以二};
}
图\ref{fig:paircounta}中两条有序记录对应同一次碰撞，因此总计数除以二。有限$N$时还要排除粒子与自身配对，将靶粒子密度改为$(N-1)/V$，同一稀薄无关联模型给
\[\mathcal R=\frac{N(N-1)\sigma}{\sqrt2V}
\sqrt{\frac{8\kb T}{\pi m}}.\]
宏观气体$N\gg1$时，它与前式的相对差为$1/N$。
\end{solution}

\begin{exercise}[小孔两侧的平衡条件]
两容器以很小的薄孔相通，孔的尺度远小于两侧平均自由程。两侧为相同气体，温度 $T_1,T_2$ 由外界维持。求零净粒子通量时 $p_1/p_2$，并说明零粒子通量是否意味着无能量传递。
\Diagram{fig:thermaltransq}{不同温度容器通过分子流小孔相连}{
 \draw(0,0)rectangle(5,1.6);\draw(2.5,0)--(2.5,.65);\draw(2.5,.95)--(2.5,1.6);\node at(1.2,.8){$p_1,T_1$};\node at(3.8,.8){$p_2,T_2$};
}
\end{exercise}
\begin{solution}
图\ref{fig:thermaltransq}两侧各自向孔贡献 $j_i=n_i\bar v_i/4$。零净粒子流要求 $n_1\sqrt{T_1}=n_2\sqrt{T_2}$。用 $p_i=n_i\kb T_i$ 消去数密度，得
\[
 \frac{p_1}{\sqrt{T_1}}=\frac{p_2}{\sqrt{T_2}},\qquad
 \frac{p_1}{p_2}=\sqrt{\frac{T_1}{T_2}}.
\]
热侧在同压下密度较小，而平均速率较大，结合后单向通量正比于$p_i/\sqrt{T_i}$。孔口处是分子流，两侧没有通过密集碰撞建立共同压强；恒温装置仍需补入或取走穿孔传递的净能量。

\Diagram{fig:thermaltransa}{粒子流相抵时能量流仍可不相抵}{
 \draw[axis](0,.5)--(4,.5)node[midway,above]{$j\times2\kb T_1$};\draw[axis](4,-.5)--(0,-.5)node[midway,below]{$j\times2\kb T_2$};
}
对单原子气体，每个逸出粒子带走 $2\kb T_i$，故图\ref{fig:thermaltransa}给出净能流 $2\kb j(T_1-T_2)$。只有温度也相同才无净热传递。这里由外界恒温，不要求两侧构成孤立平衡系统。
\end{solution}

\begin{exercise}[束流的角分布与圆锥筛选]
由薄孔逸出的分子在远场可视为余弦束。收集器只接受与孔法线夹角 $0\le\theta\le\theta_0$ 的分子。求收集到的比例和归一化角密度；若筛选只与方向有关，收集束的平均能量是否改变？
\Diagram{fig:coneq}{小孔法线周围的收集圆锥剖面}{
 \draw[line width=1.2pt](0,-1.2)--(0,-.15);\draw[line width=1.2pt](0,.15)--(0,1.2);
 \draw[axis](0,0)--(4,0)node[right]{法线};\draw[curve](0,0)--(3,1.2);\draw[curve](0,0)--(3,-1.2);
 \draw[->](.9,0)arc(0:21.8:.9);\node at(1.4,.4){$\theta_0$};\node[left]at(0,.3){孔};
}
\end{exercise}
\begin{solution}
图\ref{fig:coneq}是以法线为轴的三维圆锥的剖面。对整个方位角积分后，原束的极角密度为 $q(\theta)=2\sin\theta\cos\theta$，所以
\[
 P_c=\int_0^{\theta_0}q(\theta)\dd\theta
=[\sin^2\theta]_0^{\theta_0}=\sin^2\theta_0,
 \qquad q_c(\theta)=\frac{2\sin\theta\cos\theta}{\sin^2\theta_0}.
\]
收集比例的分母对应整个向外半球。方向密度 $\cos\theta/\pi$ 乘立体角环带 $2\pi\sin\theta\dd\theta$，才成为所用的极角密度。它的积分是 $\sin^2\theta$；不能仅以圆锥立体角占半球的比例 $1-\cos\theta_0$ 代替，因为偏斜方向的通量还受到余弦投影抑制。小角度时收集比例约为 $\theta_0^2$，缩小接受角主要减少粒子数，不额外筛选快慢。

\Diagram{fig:conea}{速度权重与方向权重可分离}{
 \node[draw]at(0,0){$g_e(v)\dd v$};\node at(2,0){$\times$};\node[draw]at(4,0){$q_c(\theta)\dd\theta$};\node at(2,-.8){方向筛选不改变速率密度};
}
原联合密度中速率与方向相乘，图\ref{fig:conea}的筛选只改变角因子的归一化。因此收集束仍有 $\langle E_k\rangle=2\kb T$。若孔道很长，透射概率可能依赖速率、壁面吸附和再发射，便不能仅凭这一角度模型推断。
\end{solution}

\begin{exercise}[重分子与轻分子的泄流分离]
容器内有两种理想气体，质量分别为 $m_a,m_b$，粒子数为 $N_a,N_b$，共同温度 $T$。在温度和容器体积保持不变的小孔泄流过程中，求逸出流的组分比和余气组分随时间的变化。
\Diagram{fig:isotopeq}{同一小孔对不同质量粒子的通量不同}{
 \draw(0,0)rectangle(3,1.6);\fill[white](2.95,.65)rectangle(3.05,.95);\node at(1.5,1.05){$N_a,m_a$};\node at(1.5,.5){$N_b,m_b$};\draw[->](3,.8)--(4.5,.8);
}
\end{exercise}
\begin{solution}
由图\ref{fig:isotopeq}，两种气体共享孔面积 $A_h$ 和体积 $V$。各自数流率为
\[
 J_i=\frac{A_hN_i}{V}\sqrt{\frac{\kb T}{2\pi m_i}}=\kappa_iN_i,
 \qquad \kappa_i=\frac{A_h}{V}\sqrt{\frac{\kb T}{2\pi m_i}}.
\]
所以瞬时逸出流的组分比为
\[
 \frac{J_a}{J_b}=\frac{N_a}{N_b}\sqrt{\frac{m_b}{m_a}}.
\]
微分方程分离为$\dd N_i/N_i=-\kappa_i\dd t$。从$N_i(0)=N_{i0}>0$积分，
\[\int_{N_{i0}}^{N_i(t)}\frac{\dd N}{N}
=-\kappa_i\int_0^t\dd s,
\qquad[\ln N]_{N_{i0}}^{N_i(t)}=-\kappa_i[s]_0^t.\]
所以$\ln[N_i(t)/N_{i0}]=-\kappa_it$，取指数给出下述各组分数目。若$m_a<m_b$，则$\kappa_a>\kappa_b$，轻组分在瞬时逸出流中富集，也在容器内消耗得更快。逸出流比为$J_a/J_b$，余气比为$N_a/N_b$，两者相差一个质量平方根因子。

\Diagram{fig:isotopea}{恒温泄流时两组分分别指数减少}{
 \draw[axis](0,0)--(5,0)node[right]{$t$};\draw[axis](0,0)--(0,2.3)node[above]{$N_i/N_{i0}$};
 \draw[curve]plot[domain=0:4.6](\x,{2*exp(-\x)});\draw[curve,dashed]plot[domain=0:4.6](\x,{2*exp(-\x/2)});\draw(-.06,2)--(.06,2);\node[left]at(0,2){$1$};\node at(3,1.6){轻组分减少较快};
}
两组分数目为$N_i=N_{i0}\ee^{-\kappa_it}$，余气比值为$(N_{a0}/N_{b0})\ee^{-(\kappa_a-\kappa_b)t}$。图\ref{fig:isotopea}实线示意较轻组分。外界持续供热使温度不变，$\kappa_i$才保持常数；若改为绝热容器，泄流系数中的$\sqrt{T(t)}$须与气体能量方程联立求解。
\end{solution}

\begin{exercise}[自由管模型的质量极限]
在例\ref{ex:piston}中，保持初始气体状态不变，分别取 $m_2\to\infty$ 与 $m_1=m_2=M_p$。求活塞实验室速度和管体速度，比较相对速度与可用于喷出活塞的机械能。
\Diagram{fig:recoillimitq}{同一膨胀能量分配给两个可运动物体}{
 \node[draw,minimum width=1.5cm]at(0,0){管 $m_2$};\node[draw]at(3,0){活塞 $m_1$};\draw[axis](-.8,-.6)--(-1.5,-.6);\draw[axis](3,-.6)--(4,-.6);\node at(1.5,.8){总机械能 $E$ 固定};
}
\end{exercise}
\begin{solution}
令 $E=\nu C_{V,m}T_0(1-2^{1-\gamma})$。图\ref{fig:recoillimitq}要求总动量为零，故 $u=\sqrt{2E/\mu_p}$，$v_1=m_2u/(m_1+m_2)$，$v_2=-m_1u/(m_1+m_2)$。
把相对速度代回两物体速度后，动能分配可不依赖具体气体参数写成
\[
 E_1=\frac{m_2}{m_1+m_2}E,\qquad
 E_2=\frac{m_1}{m_1+m_2}E.
\]
两项相加恢复总膨胀能 $E$。较轻的物体获得较大的速度和较多的动能，但两者动量大小始终相等、方向相反。固定管极限和等质量极限都可由这两个分配系数直接读出，避免把相对运动的约化质量误认成活塞在实验室系中的实际质量。

\Diagram{fig:recoillimita}{固定管与等质量反冲的速度分配}{
 \node at(0,1){$m_2\to\infty$};\draw[axis](0,0)--(1.6,0)node[right]{$v_1=\sqrt{2E/m_1}$};
 \node at(0,-1){$m_1=m_2$};\draw[axis](0,-2)--(1.4,-2)node[right]{$v_1=\sqrt{E/M_p}$};\draw[axis](0,-2)--(-1.4,-2)node[left]{$v_2=-v_1$};
}
固定管极限$\mu_p\to m_1$，管速度趋零，全部机械能进入活塞：
\[u\longrightarrow\sqrt{\frac{2E}{m_1}},\qquad
v_1\longrightarrow\sqrt{\frac{2E}{m_1}},\qquad v_2\longrightarrow0.\]
等质量时$\mu_p=M_p/2$，两物体各得$E/2$，且
\[u=2\sqrt{\frac E{M_p}},\qquad
v_1=\frac u2=\sqrt{\frac E{M_p}},\qquad
v_2=-\frac u2=-\sqrt{\frac E{M_p}}.\]
图\ref{fig:recoillimita}中相对速度增大，而活塞分得的能量减少；管体反冲既贡献了相对速度，也带走一半机械能。
\end{solution}
